%% Research in Engineering Design -- Springer Nature template (sn-jnl), author-year references.
%DIF LATEXDIFF DIFFERENCE FILE
%DIF DEL old.tex    Wed Oct  7 14:14:43 2026
%DIF ADD main.tex   Wed Oct  7 14:14:43 2026
%% Line numbers for review (option lineno); remove the option for the final version.
\documentclass[pdflatex,sn-basic,lineno]{sn-jnl}

\usepackage{graphicx}
\usepackage{multirow}
\usepackage{amsmath,amssymb,amsfonts}
\usepackage{amsthm}
\usepackage{booktabs}
\usepackage{array}
\usepackage{tikz}
\usetikzlibrary{arrows.meta,positioning,calc}
\usepackage{siunitx}
\usepackage{manyfoot}   % the class's begin-document footnote hook expects this package to be loaded
\usepackage{placeins}   % \FloatBarrier keeps floats from drifting past sections
\graphicspath{{figures/}}

\raggedbottom

%% Float placement: allow figures and tables to share pages with text and with each other.
\renewcommand{\topfraction}{0.9}
\renewcommand{\bottomfraction}{0.7}
\renewcommand{\textfraction}{0.07}
\renewcommand{\floatpagefraction}{0.8}
\setcounter{topnumber}{3}
\setcounter{totalnumber}{4}
\setlength{\textfloatsep}{10pt plus 3pt minus 2pt}
\setlength{\floatsep}{10pt plus 3pt minus 2pt}
\setlength{\intextsep}{10pt plus 3pt minus 2pt}
\setlength{\belowcaptionskip}{3pt}
%% Float pages fill from the top instead of centring their floats.
\makeatletter
\setlength{\@fptop}{0pt}
\setlength{\@fpsep}{12pt plus 2fil}
\setlength{\@fpbot}{0pt plus 1fil}
\makeatother
%% Table body size: the class's own 8 bp table font (its \footnotesize is 7 pt and its \scriptsize 9 pt).
\newcommand{\tabsize}{\fontsize{8}{9.5}\selectfont}
%% Reference list: the class sets it at normal size with entries 1em apart; the 8 bp table font with
%% hanging indents keeps the entries distinct on fewer pages.
\setlength{\bibsep}{0pt}
\makeatletter
\def\bibfont{\reset@font\fontfamily{\rmdefault}\fontsize{8}{9.5}\selectfont}
\makeatother
%% DOIs in the reference list as doi:..., linked to doi.org.
\newcommand{\doi}[1]{\href{https://doi.org/#1}{doi:\nolinkurl{#1}}}
%% Ragged-right paragraph column for text tables.
\newcolumntype{L}[1]{>{\raggedright\arraybackslash\hspace{0pt}}p{#1}}

%% Set \anontrue to produce an anonymised copy (no author block, no identifying declarations).
\newcommand{\pending}[1]{\textit{[#1 pending]}}
\newif\ifanon
\anonfalse
%DIF PREAMBLE EXTENSION ADDED BY LATEXDIFF
%DIF UNDERLINE PREAMBLE %DIF PREAMBLE
\RequirePackage[normalem]{ulem} %DIF PREAMBLE
\RequirePackage{color}\definecolor{RED}{rgb}{1,0,0}\definecolor{BLUE}{rgb}{0,0,1} %DIF PREAMBLE
\providecommand{\DIFadd}[1]{{\protect\color{blue}\uwave{#1}}} %DIF PREAMBLE
\providecommand{\DIFdel}[1]{{\protect\color{red}\sout{#1}}}                      %DIF PREAMBLE
%DIF SAFE PREAMBLE %DIF PREAMBLE
\providecommand{\DIFaddbegin}{} %DIF PREAMBLE
\providecommand{\DIFaddend}{} %DIF PREAMBLE
\providecommand{\DIFdelbegin}{} %DIF PREAMBLE
\providecommand{\DIFdelend}{} %DIF PREAMBLE
\providecommand{\DIFmodbegin}{} %DIF PREAMBLE
\providecommand{\DIFmodend}{} %DIF PREAMBLE
%DIF FLOATSAFE PREAMBLE %DIF PREAMBLE
\providecommand{\DIFaddFL}[1]{\DIFadd{#1}} %DIF PREAMBLE
\providecommand{\DIFdelFL}[1]{\DIFdel{#1}} %DIF PREAMBLE
\providecommand{\DIFaddbeginFL}{} %DIF PREAMBLE
\providecommand{\DIFaddendFL}{} %DIF PREAMBLE
\providecommand{\DIFdelbeginFL}{} %DIF PREAMBLE
\providecommand{\DIFdelendFL}{} %DIF PREAMBLE
\newcommand{\DIFscaledelfig}{0.5}
%DIF HIGHLIGHTGRAPHICS PREAMBLE %DIF PREAMBLE
\RequirePackage{settobox} %DIF PREAMBLE
\RequirePackage{letltxmacro} %DIF PREAMBLE
\newsavebox{\DIFdelgraphicsbox} %DIF PREAMBLE
\newlength{\DIFdelgraphicswidth} %DIF PREAMBLE
\newlength{\DIFdelgraphicsheight} %DIF PREAMBLE
% store original definition of \includegraphics %DIF PREAMBLE
\LetLtxMacro{\DIFOincludegraphics}{\includegraphics} %DIF PREAMBLE
\newcommand{\DIFaddincludegraphics}[2][]{{\color{blue}\fbox{\DIFOincludegraphics[#1]{#2}}}} %DIF PREAMBLE
\newcommand{\DIFdelincludegraphics}[2][]{% %DIF PREAMBLE
\sbox{\DIFdelgraphicsbox}{\DIFOincludegraphics[#1]{#2}}% %DIF PREAMBLE
\settoboxwidth{\DIFdelgraphicswidth}{\DIFdelgraphicsbox} %DIF PREAMBLE
\settoboxtotalheight{\DIFdelgraphicsheight}{\DIFdelgraphicsbox} %DIF PREAMBLE
\scalebox{\DIFscaledelfig}{% %DIF PREAMBLE
\parbox[b]{\DIFdelgraphicswidth}{\usebox{\DIFdelgraphicsbox}\\[-\baselineskip] \rule{\DIFdelgraphicswidth}{0em}}\llap{\resizebox{\DIFdelgraphicswidth}{\DIFdelgraphicsheight}{% %DIF PREAMBLE
\setlength{\unitlength}{\DIFdelgraphicswidth}% %DIF PREAMBLE
\begin{picture}(1,1)% %DIF PREAMBLE
\thicklines\linethickness{2pt} %DIF PREAMBLE
{\color[rgb]{1,0,0}\put(0,0){\framebox(1,1){}}}% %DIF PREAMBLE
{\color[rgb]{1,0,0}\put(0,0){\line( 1,1){1}}}% %DIF PREAMBLE
{\color[rgb]{1,0,0}\put(0,1){\line(1,-1){1}}}% %DIF PREAMBLE
\end{picture}% %DIF PREAMBLE
}\hspace*{3pt}}} %DIF PREAMBLE
} %DIF PREAMBLE
\LetLtxMacro{\DIFOaddbegin}{\DIFaddbegin} %DIF PREAMBLE
\LetLtxMacro{\DIFOaddend}{\DIFaddend} %DIF PREAMBLE
\LetLtxMacro{\DIFOdelbegin}{\DIFdelbegin} %DIF PREAMBLE
\LetLtxMacro{\DIFOdelend}{\DIFdelend} %DIF PREAMBLE
\DeclareRobustCommand{\DIFaddbegin}{\DIFOaddbegin \let\includegraphics\DIFaddincludegraphics} %DIF PREAMBLE
\DeclareRobustCommand{\DIFaddend}{\DIFOaddend \let\includegraphics\DIFOincludegraphics} %DIF PREAMBLE
\DeclareRobustCommand{\DIFdelbegin}{\DIFOdelbegin \let\includegraphics\DIFdelincludegraphics} %DIF PREAMBLE
\DeclareRobustCommand{\DIFdelend}{\DIFOaddend \let\includegraphics\DIFOincludegraphics} %DIF PREAMBLE
\LetLtxMacro{\DIFOaddbeginFL}{\DIFaddbeginFL} %DIF PREAMBLE
\LetLtxMacro{\DIFOaddendFL}{\DIFaddendFL} %DIF PREAMBLE
\LetLtxMacro{\DIFOdelbeginFL}{\DIFdelbeginFL} %DIF PREAMBLE
\LetLtxMacro{\DIFOdelendFL}{\DIFdelendFL} %DIF PREAMBLE
\DeclareRobustCommand{\DIFaddbeginFL}{\DIFOaddbeginFL \let\includegraphics\DIFaddincludegraphics} %DIF PREAMBLE
\DeclareRobustCommand{\DIFaddendFL}{\DIFOaddendFL \let\includegraphics\DIFOincludegraphics} %DIF PREAMBLE
\DeclareRobustCommand{\DIFdelbeginFL}{\DIFOdelbeginFL \let\includegraphics\DIFdelincludegraphics} %DIF PREAMBLE
\DeclareRobustCommand{\DIFdelendFL}{\DIFOaddendFL \let\includegraphics\DIFOincludegraphics} %DIF PREAMBLE
%DIF COLORLISTINGS PREAMBLE %DIF PREAMBLE
\RequirePackage{listings} %DIF PREAMBLE
\RequirePackage{color} %DIF PREAMBLE
\lstdefinelanguage{DIFcode}{ %DIF PREAMBLE
%DIF DIFCODE_UNDERLINE %DIF PREAMBLE
  moredelim=[il][\color{red}\sout]{\%DIF\ <\ }, %DIF PREAMBLE
  moredelim=[il][\color{blue}\uwave]{\%DIF\ >\ } %DIF PREAMBLE
} %DIF PREAMBLE
\lstdefinestyle{DIFverbatimstyle}{ %DIF PREAMBLE
	language=DIFcode, %DIF PREAMBLE
	basicstyle=\ttfamily, %DIF PREAMBLE
	columns=fullflexible, %DIF PREAMBLE
	keepspaces=true %DIF PREAMBLE
} %DIF PREAMBLE
\lstnewenvironment{DIFverbatim}{\lstset{style=DIFverbatimstyle}}{} %DIF PREAMBLE
\lstnewenvironment{DIFverbatim*}{\lstset{style=DIFverbatimstyle,showspaces=true}}{} %DIF PREAMBLE
%DIF END PREAMBLE EXTENSION ADDED BY LATEXDIFF

\begin{document}

\title[Evaluating design-stage manufacturability decisions]{Evaluating design-stage manufacturability decisions against the resolution of production}

\ifanon
\author*[1]{\fnm{Anonymised} \sur{for review}}
\affil*[1]{\orgname{Affiliation withheld for review}}
\else
\author*[1]{\fnm{H\"useyin Tayyer} \sur{Canseven}}\email{huseyin.canseven@lut.fi}

\affil*[1]{\orgdiv{Department of Electrical Engineering, LUT School of Energy Systems}, \orgname{LUT University}, \orgaddress{\street{Yliopistonkatu 34}, \city{Lappeenranta}, \postcode{53850}, \country{Finland}}}
\fi

%% BEGIN ABSTRACT
\abstract{Design-stage manufacturability decisions rest increasingly on simulation and machine learning, whose fitness to decide is rarely evaluated against production. \DIFdelbegin \DIFdel{A framework is }\DIFdelend \DIFaddbegin \DIFadd{The }\DIFaddend proposed \DIFdelbegin \DIFdel{whose }\DIFdelend \DIFaddbegin \DIFadd{framework's }\DIFaddend unit, the production floor, is the difference that drift and scatter of production and measurement produce between two batches of one design within a run. For each relation of a new design to produced evidence, a rule returning meets, fails or trial has a decisive distance in floors, beyond which at least 95\,\% of its verdicts are correct, and a safe distance, beyond which at most 5\,\% are wrong; a procedure adds guard bands for single verdicts. Fourteen rules were evaluated at 95\,\% conformance on open data of 18 deep-drawn alternatives of 500 parts, mostly process settings, with matched simulations \DIFdelbegin \DIFdel{; }\DIFdelend \DIFaddbegin \DIFadd{(}\DIFaddend a draw-in floor\DIFdelbegin \DIFdel{is }\DIFdelend \DIFaddbegin \DIFadd{: }\DIFaddend 0.05\,mm\DIFaddbegin \DIFadd{)}\DIFaddend . With a produced sibling, the nearest produced setting was right for requirements over 3.4 floors from the truth (\DIFdelbegin \DIFdel{0.85 for }\DIFdelend near-replicates \DIFaddbegin \DIFadd{0.85}\DIFaddend , \DIFdelbegin \DIFdel{4.9 for }\DIFdelend resolvable siblings \DIFaddbegin \DIFadd{4.9}\DIFaddend ). For a new setting the best rules needed 4.7 floors interpolating, 4.9 extrapolating to 500\,kN and 8.9 to 100\,kN, with a stroke-speed change; for a new family the best was safe at 16--17 floors in either family's floor. \DIFdelbegin \DIFdel{Release-level requirements }\DIFdelend \DIFaddbegin \DIFadd{Requirements at a performance index of 1.33 }\DIFaddend lay inside every decisive distance except a near-replicate's. With six to nine produced alternatives at two or three forces, the simulation as used, offset or rescaled, helped only across families\DIFdelbegin \DIFdel{, and }\DIFdelend \DIFaddbegin \DIFadd{; }\DIFaddend learned 90\,\% intervals covered 85--96\,\% with a sibling, under 50\,\% extrapolating or across families. \DIFdelbegin \DIFdel{The framework }\DIFdelend \DIFaddbegin \DIFadd{This }\DIFaddend gives AI-based manufacturability support a production-referenced evaluation standard.}
%% END ABSTRACT

\keywords{design for manufacturing, machine learning, design decision making, verification and validation, operating characteristic, deep drawing}

\maketitle

%% ----------------------------------------------------------------------
\section{Introduction}\label{sec:intro}

Whether a design alternative can be made to its requirements is decided long before it is made. In sheet-metal forming the decision rests on a finite-element simulation \citep{banabic2010sheet}: the team compares the predicted characteristics with the requirements and releases the alternative, changes it or sends it to a physical trial. Simulation has moved such decisions forward \citep{thomke1998simulation}, and machine learning makes the loop faster: surrogates replace simulations, models learn manufacturability from earlier designs, and checks are automated \citep{yan2026machine, ghadai2018learning, zhang2018featurenet}. Neither has changed the question the team has to answer: \emph{can this prediction be relied on to decide?} Whitney argued in this journal that mechanical design cannot be verified by simulation the way integrated circuits are \citep{whitney1996why}, and production variation, with which development organisations still contend \citep{thornton2000more}, adds a further reason: whether a simulation-based decision is credible is a property of production, and production evidence is scarce at the design stage.

This paper treats the design-stage manufacturability decision itself as the object of study. How can the fitness of a design-stage prediction to decide be expressed on a common production-referenced scale, for any predictor and characteristic (RQ1)? How does it depend on the amount of production evidence and on how the new design relates to what has been produced (RQ2)? And what do simulation and machine learning contribute once production evidence exists (RQ3)? Decision-based design \citep{hazelrigg1998framework}, design margins \citep{eckert2019design} and the planning of verification and testing \citep{thomke2001sequential, salado2018mathematical, tahera2019testing} all need to know how far a prediction can be relied on.

The framework proposed here has the \emph{production floor} as its unit: the difference that drift and scatter \DIFaddbegin \DIFadd{of production and measurement }\DIFaddend produce between two batches of one design within a run, a floor in the sense of a noise floor. A rule that turns the design-stage evidence into a verdict (\DIFdelbegin \DIFdel{the alternative }\DIFdelend \emph{meets}\DIFdelbegin \DIFdel{the requirement}\DIFdelend , \emph{fails} \DIFdelbegin \DIFdel{it, or goes to }\DIFdelend \DIFaddbegin \DIFadd{or }\DIFaddend a physical \emph{trial}) is characterised, for each \emph{evidence relation} between a new design and the produced alternatives, by two points on its operating characteristics in floors, its \emph{decisive} and \emph{safe distances}, and a procedure turns them into a guard band for single verdicts. The decisions in scope are those of a team that has produced variants of a design family, in process planning and tryout and when production knowledge is carried to a new family.

The evaluation uses open data: 9,000 deep-drawn and cut cups of dual-phase steel, produced as 18 design--process alternatives of 500 consecutive parts each \citep{baum2026rddac}, with finite-element simulations of the same process \citep{baum2025ddacs}. Within each of two cup geometries the alternatives differ in blank-holder force and lubrication on one tool, so two of the three relations concern process settings, and the geometries give two transfers between families. Each alternative is judged as a new design by fourteen rules at a 95\,\% conformance level, against what its 500 parts did. In the terms of the design research methodology \citep{blessing2009drm}, this is a prescriptive study with an initial support evaluation: it shows whether the framework discriminates between rules and evidence situations, not whether design teams decide better with it, and its findings on one dataset are hypotheses for confirmation. After the positioning (Section~\ref{sec:background}), the contributions are the framework (Section~\ref{sec:method}), an evaluation design that separates the evidence relations and isolates what the simulation adds (Section~\ref{sec:demo}), the empirical answers for deep drawing (Section~\ref{sec:results}), and guidance and a reporting standard for evaluating simulation- and AI-based manufacturability support (Section~\ref{sec:discussion}).

%% ----------------------------------------------------------------------
\section{Background and related work}\label{sec:background}

\subsection{Manufacturability evaluation and machine learning}

Design for manufacturing (DFM) codifies what a process can make into rules, guidelines and cost models \citep{boothroyd2010product}, and quantitative manufacturability models have driven design decisions in this journal \citep{latrobe2003design}. Machine learning learns manufacturability from CAD models \citep{ghadai2018learning, zhang2018featurenet}, from process data and simulated design spaces \citep{yan2026machine}, and generative models propose designs directly \citep{regenwetter2022deep}. Learned manufacturability models are evaluated by held-out accuracy on labelled designs, and language and vision-language models by case studies \citep{makatura2024how} and benchmark questions \citep{picard2025concept}, not by the decisions they support against the production of the alternatives they are asked about.

\subsection{Design decisions, design types and production evidence}

Decision-based design casts design as a choice between alternatives with uncertain outcomes \citep{hazelrigg1998framework}. How much a team knows about a new alternative depends on the kind of design, variant, adaptive or original \citep{pahl2007engineering}, organised across products by platform design \citep{simpson2006product}; the evidence relations of this paper are their counterparts for production evidence, the design type concerning the novelty of the solution and the evidence relation the novelty of its realisation relative to what has been produced. Production knowledge enters design through process-capability data \citep{thornton2004variation, tata1999process}. Unresolved uncertainty is absorbed by margins \citep{eckert2019design}, whose part that covers uncertainty is distinct from the part that absorbs change \citep{eckert2017safety}. Tests and prototypes, the alternative to a margin, are modelled as a trade between the costs of tests and of errors \citep{thomke1998simulation, thomke2001sequential, salado2018mathematical} and studied as design activities \citep{tahera2019testing}. When upstream information is fit to act on is the question of overlapped development \citep{krishnan1997model, terwiesch2002exchanging}, and set-based concurrent engineering codifies production knowledge as feasibility limits \citep{sobek1999toyota}. All of these assume a measure of how far a design-stage prediction can be relied on; the framework supplies one in a unit set by production.

\subsection{Simulation credibility, calibration and validation}

Verification and validation ask whether a model is accurate enough for its intended use; the distinction between its validation and application domains \citep{oberkampf2010verification} is the question that the evidence relation asks of a design, and risk-informed credibility assessment scales the required evidence with the consequence of a wrong decision \citep{asme2018assessing}. Bayesian calibration adds a discrepancy model and calibration parameters to a simulator \citep{kennedy2001bayesian}, multi-fidelity models scale a cheap model to an expensive one \citep{kennedy2000predicting}, and a model may be relied on only within its valid input domain \citep{malak2010using}. Conformal prediction gives distribution-free intervals under exchangeability \citep{lei2018distribution, barber2021predictive}, also under covariate shift \citep{tibshirani2019conformal} and for decision risks \citep{angelopoulos2024conformal}. In sheet-metal forming, friction \citep{hol2012advanced} and the unloading behaviour of the material \citep{yoshida2002model} are main sources of simulation error, simulation datasets of deep drawing support surrogate and transfer learning \citep{baum2025ddacs, heinzelmann2025comprehensive}, and the deviation between simulation and parts has been decomposed statistically \citep{baum2026statistical}. How close to a requirement such intervals and accuracy statements can be relied on to decide, and how much production evidence that takes, remains open.

\subsection{Decisions with abstention and their operating characteristics}

In conformity assessment a measured value proves conformance, proves non-conformance or leaves the decision open, with guard bands set by the measurement uncertainty and the risks of false acceptance and rejection as the measures of a decision rule \citep{iso2017decision}. In selection procedures and acceptance sampling the operating characteristic gives the probability of a decision against the true quality, with an indifference zone inside which no decision is required \citep{bechhofer1954single}. A classifier with a reject option trades errors against abstentions \citep{chow1970optimum, geifman2017selective}, learning to defer passes difficult cases to an expert \citep{madras2018predict}, and appropriate reliance on automation requires that its reliability be visible to its users \citep{lee2004trust}, also in design teams working with AI \citep{zhang2021cautionary}. The decisive and safe distances are operating characteristics of this kind, and the procedure's guard band is the design-stage counterpart of a conformity-assessment one.

\subsection{Production variation and measurement}

Statistical quality control compares short- and long-term variation with tolerances through capability and performance indices \citep{montgomery2019introduction}, with release criteria such as a performance index of at least 1.33 \citep{aiag2006production} and time-dependent models for drifting processes \citep{iso2026process}. Precision experiments define the repeatability limit, the difference that two results under repeatability conditions exceed with a probability of 5\,\% \citep{iso2025repeatability}, and intermediate precision under conditions that differ in time, operator or equipment \citep{iso2023intermediate}; measurement systems analysis separates repeatability from stability \citep{aiag2010measurement}. The production floor is a precision limit of this kind for batch centres within a run (Section~\ref{sec:floor}); what is new is its use not to judge whether two results agree but as the unit in which a design decision rule is scored and its guard band set. Table~\ref{tab:positioning} places the framework among them: it does not replace calibration or conformal prediction but scores them as decision rules against production in a common unit.

\begin{table}[t]
\tabsize
\setlength{\tabcolsep}{3.5pt}
\caption{Positioning among methods that support or evaluate design-stage manufacturability decisions}\label{tab:positioning}
\begin{tabular}{@{}L{3.0cm}L{2.5cm}L{3.0cm}L{1.9cm}L{1.6cm}@{}}
\toprule
Approach & Evidence & Output & Scored against production & Abstains \\
\midrule
DFM rules and cost models \citep{boothroyd2010product} & codified process knowledge & guideline compliance, cost & no & no \\
Design margins \citep{eckert2019design} & experience, analysis & buffer to the requirement & rarely & no \\
V\&V, credibility assessment \citep{oberkampf2010verification, asme2018assessing} & physics model, validation tests & accuracy in physical units; required evidence & on test cases & can require tests \\
Bayesian calibration \citep{kennedy2001bayesian} & simulations, experiments & predictive distribution & not as decisions & implicitly \\
Conformal prediction \citep{lei2018distribution} & exchangeable calibration data & interval with coverage & coverage; per decision possible & yes \\
Conformity assessment, OC curves \citep{iso2017decision, bechhofer1954single} & measurement and its uncertainty & accept, reject, undecided; risks & on the measured item & yes \\
Selective classification \citep{chow1970optimum, geifman2017selective} & labelled data & class or reject; risk--coverage & on held-out data & yes \\
ML manufacturability models \citep{ghadai2018learning, yan2026machine} & historical or simulated designs & class or value & on held-out data & rarely \\
This paper & simulations, produced alternatives & verdict; decisive and safe distance in floors by evidence relation & yes, per verdict & yes \\
\botrule
\end{tabular}
\end{table}

%% ----------------------------------------------------------------------
\section{The framework}\label{sec:method}

\subsection{The decision problem and its scope}\label{sec:problem}

A design alternative $a$ is a combination of product geometry and process settings produced as a series of parts, over which a quality characteristic $c$ varies. A requirement is an upper limit $R$ on the characteristic, or on its absolute value for characteristics whose target is zero, that a share $p$ of the parts must meet: alternative $a$ is \emph{adequate} for $c$ when the $p$-quantile of its parts satisfies $q_p(a,c)\le R$. Lower or two-sided limits follow by a change of sign or by the absolute deviation from the target. At the design stage $q_p$ is unknown. A decision rule $r$ maps the available evidence to an interval $[L_r, U_r]$ for $q_p$ and to a verdict
\begin{equation}
V_r = \begin{cases} \text{meets} & U_r \le R,\\ \text{fails} & L_r > R,\\ \text{trial} & \text{otherwise.} \end{cases}
\label{eq:verdict}
\end{equation}
The three verdicts mirror the decision rules of conformity assessment \citep{iso2017decision}, here for the uncertainty of a design-stage prediction. A rule that returns a point ($L_r=U_r$) always decides. The evidence consists of simulations of $a$ and of the production record of other alternatives, the \emph{calibration alternatives}. The demonstration decides on $p=0.95$, which for a normal characteristic corresponds to a performance index of about 0.55. Release criteria of series production, a performance index of 1.33 or more \citep{aiag2006production}, concern nonconforming fractions in parts per million that a series of 500 parts cannot estimate without a parametric tail; Section~\ref{sec:res-requirements} locates such limits relative to the distances, without deciding them.

\subsection{Production floor, effect-to-scatter ratio and minimum resolvable change}\label{sec:floor}

Production evidence comes as series of consecutive parts, each cut into batches of $B$ consecutive parts, and each batch is summarised by a robust centre $m$ (here the 10\,\% trimmed mean). The production floor of characteristic $c$ is
\begin{equation}
F_c = Q_\alpha\bigl\{\,|m_i - m_j| : i< j \text{ batches of the same alternative}\,\bigr\},
\label{eq:floor}
\end{equation}
the $\alpha$-quantile of the difference between two batches of one design. Under a stationary normal model with a standard deviation $\sigma_\mathrm{b}$ between batch means and $\sigma_\mathrm{w}$ between parts, $F_c = 1.96\sqrt{2}\sqrt{\sigma_\mathrm{b}^2+\sigma_\mathrm{w}^2/B}$ at $\alpha=0.95$, the form of the repeatability limit of precision experiments \citep{iso2025repeatability}. Because it pools batch pairs up to 450 parts apart and grows with their lag (Section~\ref{sec:res-floor}), it is closer to a time-different intermediate-precision limit \citep{iso2023intermediate} of batch centres within a run; in the terms of ISO 22514-2 the runs follow a time-dependent model with varying location \citep{iso2026process}. The reference protocol of this paper is $B=50$, ten batches of a 500-part run, all pairs of batches and $\alpha=0.95$, with per-geometry floors where geometries differ. The floor rests on three assumptions: batches of one design differ only by drift and scatter within a run (A1); the variation within the runs from which it is estimated is like that within future runs (A2); and the measuring system is repeatable and its measurement resolution small against the floor (A3, checked in Section~\ref{sec:res-floor}).

Four properties matter for its use. It has the unit of the characteristic, so distances in floors are on a common scale for all characteristics and predictors, though equal distances do not mean equal nonconforming fractions (Table~\ref{tab:floors}). It is a property of a protocol: under drift it grows with the time between the batches compared and shrinks with $B$, and $\sigma_\mathrm{b}$ and $\sigma_\mathrm{w}$ are reported so that it can be recomputed. It contains whatever drifts between batches, process and measuring system alike, so without a reference artefact scanned through the series it is a production-and-measurement floor. And it excludes the variation between runs, which makes it a lower bound of the variation of series production. A new family, whose floor is unknown before it is produced, is scored in the floor of the produced family, and because the floor resolves batch centres, the reproducibility of the decided quantile is reported alongside.

Two alternatives $a$ and $b$ are resolvable for $c$ when their effect-to-scatter ratio
\begin{equation}
\mathrm{ESR}_c(a,b) = |\bar m_a - \bar m_b| \,/\, F_c
\label{eq:esr}
\end{equation}
is at least one, $\bar m$ being the centre of all parts of an alternative; this is a practical threshold, not a test. For a continuous design or process parameter $x$ with a local sensitivity $s_c = \partial \bar m/\partial x$, the minimum resolvable change is $F_c/|s_c|$: a change of $x$ by less moves the centre of the characteristic by less than two batches of one design differ. It is local and concerns centres; a change below it can still move the tail.

\subsection{Decision rules}\label{sec:rules}

Fourteen rules span what a design team can do with its evidence (Table~\ref{tab:rules}); the first seven and the force trend M1sn carry the main results. Beside the nominal simulation $s_0$ (M0), the bias-corrected simulation (M1) adds the median offset between the calibration alternatives' $q_p$ and their $s_0$; the scaled simulation (M1s) fits $q_p = a + b\,s_0$ to them by least squares, a linear multi-fidelity correction \citep{kennedy2000predicting}\DIFdelbegin \DIFdel{that rescales a simulated trend of the wrong size}\DIFdelend . The envelope rule (M2) adds the median offset and the largest absolute residual of the calibration alternatives to the upper end of the simulation envelope over plausible incoming conditions, which vary sheet thickness and friction but not material properties. M5 models the offset between $q_p$ and the envelope as a Gaussian process \citep{rasmussen2006gaussian} with a 90\,\% predictive interval, whose noise variance is bounded below by the variance between the calibration alternatives' series of one geometry and force, so that a new series of a produced setting is not predicted more precisely than produced series reproduce each other; M5n is the same process of $q_p$ without the simulation. The nearest produced setting (NN) is a capability-database look-up that uses neither simulation nor learning. The seven further rules include a friction-calibrated simulation \citep{kennedy2001bayesian} (Mc), part-level gradient boosting \citep{friedman2001greedy} with a jackknife+ interval \citep{barber2021predictive}, with and without the simulation (M3, M3n), and the least-squares line of $q_p$ in force (M1sn). The nominal simulation and the envelope do not depend on the lubrication pattern, so within a geometry they are functions of the force alone: M1sn is the counterpart of M1s without the simulation, and a lubrication variant gives M0, M1, M1s, M2 and M5 no simulated information that distinguishes it from its siblings.

With $n$ calibration alternatives, the largest-residual margin covers a new exchangeable residual with probability $n/(n+1)$ for a fixed offset and at least $(n-1)/(n+1)$ for an estimated one; it is the split-conformal quantile \citep{lei2018distribution} at the highest attainable level, not a 90\,\% interval. A distribution-free interval cannot certify a miss rate below $1/(n+1)$, so a miss rate of 5\,\% needs at least 19 produced alternatives in the relevant relation. With six to nine calibration alternatives the jackknife+ bounds at the level of 0.9 fall at or beyond the extreme leave-one-out values, which are used instead; under exchangeability these guarantee a coverage of $1-2/(n+1)$, 0.71 to 0.80 \citep{barber2021predictive}, approximately because the residuals are centred on their median. The Gaussian-process interval is model-based. The level trades decisiveness for safety, as the cost of errors would set the threshold of a reject option \citep{chow1970optimum}; 0.9 is used here, 0.8 and 0.95 in the robustness analysis.

\begin{table}[t]
\tabsize
\setlength{\tabcolsep}{3.5pt}
\caption{The decision rules; the first seven and M1sn carry the main results}\label{tab:rules}
\begin{tabular}{@{}L{3.5cm}L{4.9cm}L{1.9cm}L{1.4cm}@{}}
\toprule
Rule & Prediction of $q_p$ & Interval & Evidence \\
\midrule
M0 nominal simulation & simulation at nominal thickness and friction & point & none \\
M1 bias-corrected simulation & M0 plus the median offset & point & $q_p$ \\
M1s scaled simulation & $a+b\,$M0, least squares & point & $q_p$ \\
M2 envelope + residual margin & envelope top plus the median offset & largest residual & $q_p$ \\
M5 GP-corrected envelope & envelope plus a GP offset over the descriptors & GP, 90\,\% & $q_p$ \\
M5n GP of $q_p$ & GP of $q_p$ over the descriptors & GP, 90\,\% & $q_p$ \\
NN nearest produced setting & mean $q_p$ at the nearest produced setting & point & $q_p$ \\
\midrule
M0w process-window worst case & simulated range over the incoming conditions & range & none \\
Mc friction-calibrated simulation & friction calibrated on the draw-in, plus offset & point & centres, $q_p$ \\
M3 simulation + boosting & matched simulation plus boosted discrepancy & jackknife+ & parts \\
M3n boosting, no simulation & boosted model of the characteristic & jackknife+ & parts \\
M1n median of the produced & median $q_p$ of the calibration alternatives & point & $q_p$ \\
M1sn force trend & $a+b\,$force, least squares & point & $q_p$ \\
M2n median + residual margin & M1n with the largest-residual margin & largest residual & $q_p$ \\
\botrule
\end{tabular}
\footnotetext{Fitted on the calibration alternatives. Descriptors: geometry indicator, blank-holder force and lubrication rank (the order of the oil amount). GP: Gaussian process. M3 and M3n at part level; M3n is named M4 in the released code. M6, a conformalised GP, is a robustness variant (Table~\ref{tab:distances}).}
\end{table}

\subsection{Operating characteristics: decisive and safe distances}\label{sec:distances}

A rule is scored by judging each produced alternative as a new design at requirements placed at a known distance from the truth,
\begin{equation}
R = q_p(a,c) + d\,F_c ,
\label{eq:requirement}
\end{equation}
so that $a$ is adequate exactly when $d\ge 0$. For a case $(a,c)$ with interval $[L,U]$, let $u=(U-q_p)/F_c$ and $v=(q_p-L)/F_c$ be the distances by which the interval excludes the truth from above and below. At $d=+\delta$ the verdict is correct when $u\le\delta$ and a false reject when $-v>\delta$; at $d=-\delta$ it is correct when $v<\delta$ and a false accept when $-u\ge\delta$. Let $\kappa(\delta)$ and $\omega(\delta)$ be the shares of correct and of wrong verdicts among all feasible verdicts at $|d|=\delta$: every case gives one verdict on an adequate alternative and, where $q_p-\delta F_c\ge0$, one on a failing alternative, and every verdict has equal weight. As $\delta$ grows, failing-side requirements become infeasible for characteristics close to zero, so the failing side carries half the weight up to a few floors and a third at about 100 floors (98 of 126 failing-side requirements remain feasible at 16 floors and 63 at 108). The \emph{decisive distance} $\delta_\mathrm{d}$ of a rule is the smallest $\delta$ beyond which $\kappa\ge 0.95$, and its \emph{safe distance} $\delta_\mathrm{s}$ the smallest $\delta$ beyond which $\omega\le 0.05$.

The distances are exact up to the 0.05-floor grid on which they are evaluated: $\kappa$ and $\omega$ change only where $\delta$ crosses one of the case's $u$, $v$, $-u$ or $-v$, so the distances are order statistics of these exclusion distances; for a rule that always decides, $\delta_\mathrm{d}=\delta_\mathrm{s}$ is close to the 90\,\% quantile of its absolute error in floors. Always $\delta_\mathrm{s}\le\delta_\mathrm{d}$, and between the two the rule sends designs to trial rather than misleading; one-sided distances for failing alternatives (false accepts) and adequate ones (false rejects) matter when the errors cost differently. The distances are points on operating characteristics over a population of design situations: $\delta_\mathrm{d}$ is the analogue of the indifference zone of a selection procedure \citep{bechhofer1954single} at a probability of correct selection of 0.95. A trial is not error-free: a short run reproduces the 95th percentile only to about one floor within a run (Section~\ref{sec:res-floor}) and, being a run, also carries the variation between runs, so no decision is reliable closer to the requirement than the trial replacing it.

\subsection{Evidence relation, calibration budget and applicability guard}\label{sec:budget}

How a new design relates to the produced alternatives is part of the decision situation. With a \emph{sibling}, an alternative at the same process setting has been produced that differs only in the factor varied at a fixed setting (here the lubrication pattern), which production may not resolve; a \emph{new process setting} lies between produced settings (interpolation) or outside them (extrapolation); a \emph{new family} shares no produced alternative. The calibration budget is the dependence of the distances on the number $k$ of produced alternatives, over all subsets of size $k$ classified by relation. An applicability guard \DIFdelbegin \DIFdel{returns }\emph{\DIFdel{trial}} %DIFAUXCMD
\DIFdel{when the new design }\DIFdelend \DIFaddbegin \DIFadd{flags a new design that }\DIFaddend lies outside the calibration: its geometry is not calibrated (family guard), a descriptor lies outside the calibrated range (range guard), or its simulated $q_p$ lies outside the calibration alternatives' simulated values (novelty guard), in the spirit of a valid input domain \citep{malak2010using}.

\subsection{Procedure}\label{sec:procedure}

Figure~\ref{fig:procedure} turns the framework into six steps: steps 1 to 5 once per design family, repeated when production evidence is added, and step 6 for every new design. Steps 1 to 4 frame the decision (family, alternatives, characteristics, form of the requirement, conformance level), measure the floor with a stated protocol, state which differences production resolves and calibrate the candidate rules with and without the simulation. Step 5 qualifies them by their distances for each evidence relation and fixes for each rule and relation a \emph{guard band} $g$: the smallest observable margin between the predicted interval and the requirement beyond which at least 95\,\% of the decided verdicts were correct. Being conditional on what is observed, $g$ depends on the distribution of requirement distances that the team expects, not on rule and relation alone. The reference prior weights the 47 distances of the \DIFdelbegin \DIFdel{evaluation }\DIFdelend \DIFaddbegin \DIFadd{requirement }\DIFaddend grid within $\pm20$ floors equally (0, $\pm0.5$ to $\pm10$ in steps of 0.5, $\pm12$, $\pm15$ and $\pm20$; 87\,\% of the weight within $\pm10$), drops infeasible requirements and searches $g$ in steps of 0.25 floors up to 20\DIFdelbegin \DIFdel{; it represents }\DIFdelend \DIFaddbegin \DIFadd{, beyond which the rules send most designs within the prior's range to trial; it is a reference for }\DIFaddend requirements set from a few to some tens of floors from the true 95th percentile, denser where verdicts are hardest\DIFdelbegin \DIFdel{. If the distances are too large for the requirements of coming designs, another variant is produced}\DIFdelend \DIFaddbegin \DIFadd{, which a team replaces by the distribution of its own requirements}\DIFaddend . Step 6 establishes the relation of the new design, applies the rule qualified for it\DIFdelbegin \DIFdel{and the guards}\DIFdelend , widens the rule's interval by $g$ floors and returns meets or fails with the margin in floors, or a trial when \DIFdelbegin \DIFdel{a guard fires, }\DIFdelend no rule is qualified or the requirement lies inside the widened interval\DIFaddbegin \DIFadd{. The guard band of the relation takes the place of the range and family guards, which here coincide with the extrapolated and new-family relations; the novelty guard is reported as a diagnostic, not applied, since most of its refusals would have been right}\DIFaddend ; Section~\ref{sec:res-procedure} reports $g$ under other priors and works one decision through.

\begin{figure}[t]
\centering
\input{fig_procedure}
\caption{The procedure: steps 1--5 once per design family, step 6 for every new design; dashed boxes are the evidence it uses}\label{fig:procedure}
\end{figure}

%% ----------------------------------------------------------------------
\section{Demonstration on deep drawing}\label{sec:demo}

\subsection{Alternatives and production data}\label{sec:data}

The production data are the Real Deep Drawing and Cutting dataset (RDDAC) of the University of Stuttgart \citep{baum2026rddac, baum2026statistical}. Square blanks of DP600 steel, 210\,mm wide and about 1\,mm thick, were deep drawn into cups (OP10) and cut (OP20) on one tool set with two geometries, concave and convex side walls, which differ in base curvature, wall angle and bottom radius at once. Each geometry was run at three blank-holder forces (100, 300 and 500\,kN) and three lubrication patterns (coarse, medium and fine, with series medians of 1.15--1.48, 0.92--1.23 and 0.85--1.13\,g/m$^2$ of oil): 18 alternatives, each produced as one series of 500 consecutive parts, the change of geometry being the only change of product design. Every part has press signals, sheet and oil-film traverses and 3D scans after both operations. The series warm up by 1.0--5.0\,K over their first 150 parts from punch temperatures of 20.5--29.6\,$^\circ$C, their median sheet thickness is 0.986--0.993\,mm, and the 100\,kN series ran at 47\,mm/s against 78--79\,mm/s. The documentation states neither the order and dates of the series nor whether the parts were scanned in production order.

\subsection{Simulations}\label{sec:sims}

The simulations come from the Deep Drawing and Cutting Simulations dataset (DDACS) of the same group \citep{baum2025ddacs, heinzelmann2025comprehensive}: LS-DYNA quarter models with nodes about 1\,mm apart, a scaled DP600 flow curve, constant Coulomb friction and springback after both operations; the yield criterion and unloading model, which govern springback \citep{yoshida2002model}, are not documented, and the stroke speed is not represented. For the RDDAC tools, 396 simulations vary the sheet thickness (0.95--1.00\,mm) and the friction coefficient (0.05--0.15) at each geometry and force. A part is matched to the simulation with the nearest thickness and the friction coefficient that the dataset's linear mapping assigns to its oil film, an assumption of the dataset's authors on which Mc, M3 and every matched simulation depend. The nominal simulation of an alternative is at 0.99\,mm, the grid value nearest the median thickness, and friction 0.10.

\subsection{Quality characteristics and their measurement}\label{sec:qcs}

Seven characteristics were extracted with the same definitions from the scans and the simulated nodes, referred to the plane of the cup bottom: the mid-side and corner draw-in and the flange waviness, process indicators on the flange that the cut removes, and four characteristics a drawing would carry: the wall angle (design angle 20$^\circ$ concave, 10$^\circ$ convex), the cup depth, and the arm angle and bottom dome of the cut part, the last two with requirements on absolute values. The four sides of a part measure several of them twice, which exposes the measurement: the flange width along the scan lines scatters 4 to 18 times more than across them, so the mid-side draw-in is taken across the scan lines, and the wall angle is the mean of the east and west walls. The mid-side draw-in is quantised in steps of 0.019\,mm (0.38 floors) by the pixel pitch. Differences between scans and simulations, such as the scanned surface against the shell mid-surface, enter the offsets between the two but not the floors.

\subsection{Evaluation design}\label{sec:evaldesign}

Every produced alternative is judged as a new design under grouped cross-validation \citep{roberts2017cross}, one scheme per evidence relation. \emph{New variant}: the alternative alone is held out and the other eight of its geometry calibrate, two siblings at its force among them. \emph{New process setting}: the three alternatives of one force of a geometry are held out and the six others calibrate; 300\,kN is interpolated, 100 and 500\,kN extrapolated. \emph{New family}: one geometry calibrates the rules for the other (two transfers). Pooled versions of the first two add the other geometry's alternatives to the calibration set. Every calibrated rule is evaluated with and without the simulation (M1/M1n, M1s/M1sn, M2/M2n, M3/M3n, M5/M5n), isolating what the simulation contributes, and the nearest produced setting is the baseline without simulation and learning. A scope comprises 126 cases (18 alternatives, 7 characteristics).

The intervals of the distances come from 1000 resamples of the held-out alternatives within each geometry, each with a block-bootstrap replicate of every true $q_p$, common to all rules so that differences get paired intervals; they hold the fitted rules fixed. A second bootstrap refits every alternative-level rule after drawing the three lubrication patterns of each geometry with replacement, at least two of them distinct, which keeps every case in its evidence relation (Table~\ref{tab:robust}). Its 200 resamples contain 48 of the 49 distinct configurations; in three of four draws a pattern is duplicated, so that a new variant keeps one distinct sibling and five distinct calibration alternatives, and the duplicated series enter the noise floor of the Gaussian processes as identical pairs. Its intervals thus show the effect of losing a pattern, not \DIFdelbegin \DIFdel{the sampling uncertaintyof the reported values}\DIFdelend \DIFaddbegin \DIFadd{sampling uncertainty}\DIFaddend .

\subsection{Computation and development record}\label{sec:computation}

The constants of the floor and the requirement grid were fixed before the decision results were seen, as the repository history records. The rules were not: M0--M5 were defined first; M0w, Mc, M1n, NN, M2n, M5n, the grouped schemes, the exact estimator and the noise floor of the Gaussian processes were added after the first results; a review of that version added M1s, the jackknife+ intervals, the between-run noise floor and the analyses of Sections~\ref{sec:res-drivers} and \ref{sec:res-procedure}; and a second review added M1sn, the guard bands under other priors and the new family in the produced family's floor. The comparison of rules is therefore exploratory and carries the optimism of a selection made on the evaluation data \citep{cawley2010over}; every variant tried is reported in the paper or in the archived results \citep{canseven2026archive}, which Table~\ref{tab:mapping} maps to the text.

%% ----------------------------------------------------------------------
\section{Results}\label{sec:results}

\subsection{What the production floor contains}\label{sec:res-floor}

In 95\,\% of batch pairs, two batches of 50 consecutive parts of one alternative differ by up to 0.051\,mm in mid-side draw-in and 0.078$^\circ$ in wall angle (Table~\ref{tab:floors}); one floor is 0.7 to 1.7 long-term standard deviations of the parts. The floors are set by drift within the run: without the time order they shrink by a factor of 2.1 to 4.5, and batches one position apart differ by 0.51--0.80 floors and nine apart by 1.1--1.5 (averages over the two geometries). The normal model of Section~\ref{sec:floor} reproduces the floors to within 12\,\%; for subgroups of five parts, as sampled inspection would measure, it gives floors 1.1 to 2.1 times as large (median 1.3). Independent part-to-part noise, scanner repeatability included, would produce 15--33\,\% of each floor, and the measurement resolution is small against the floor except for the mid-side draw-in; A3 holds for repeatability and \DIFaddbegin \DIFadd{measurement }\DIFaddend resolution, but the stability of the scanner cannot be checked without a reference artefact. Nor do the process signals identify the source of the drift: batch centres of press force, punch temperature, sheet thickness and oil film predict it out of fold with $R^2$ of $-0.23$ to $0.20$. Variation between runs is not in the floor: the three lubrication series of one geometry and force differ by up to 1.2 (waviness) to 9.8 floors (dome), and a floor taken across them is 1.5 to 9.1 times the reference. Within a run the decided statistic reproduces to about one floor (the 95th percentiles of two batches of 100 parts differ by 0.92--1.49 floors in 95\,\% of pairs), so distances below about one floor are beyond what production can confirm, and a trial run, which also carries the variation between runs, confirms less still.

%% BEGIN GENERATED floors
\begin{table}[t]
\tabsize
\setlength{\tabcolsep}{2.0pt}
\caption{Production floor per characteristic, pooled and per geometry, with its components}\label{tab:floors}
\begin{tabular}{@{}lrrrrrrrr@{}}
\toprule
Characteristic & $F$ & 95\,\% interval & Concave & Convex & $F/\sigma_\mathrm{LT}$ & $\sigma_\mathrm{b}/F$ & $\sigma_\mathrm{w}/F$ & Drift \\
\midrule
Draw-in, mid-side (mm) & 0.051 & 0.045--0.061 & 0.050 & 0.052 & 1.0/1.4 & 0.32/0.34 & 0.91/0.66 & 3.3 \\
Draw-in, corner (mm) & 0.046 & 0.035--0.052 & 0.053 & 0.030 & 1.4/1.2 & 0.31/0.35 & 0.65/0.77 & 4.5 \\
Flange waviness (mm) & 0.0028 & 0.0024--0.0031 & 0.0037 & 0.0021 & 1.6/0.7 & 0.33/0.30 & 0.55/1.31 & 2.5 \\
Wall angle ($^\circ$) & 0.078 & 0.066--0.10 & 0.081 & 0.078 & 1.3/0.7 & 0.32/0.26 & 0.68/1.39 & 2.1 \\
Arm angle, cut ($^\circ$) & 0.11 & 0.094--0.12 & 0.12 & 0.100 & 1.7/0.8 & 0.36/0.31 & 0.46/1.22 & 2.6 \\
Cup depth (mm) & 0.011 & 0.0096--0.012 & 0.0087 & 0.012 & 1.4/1.6 & 0.31/0.34 & 0.65/0.52 & 3.9 \\
Bottom dome, cut (mm) & 0.0070 & 0.0061--0.0081 & 0.0034 & 0.0082 & 1.2/1.6 & 0.33/0.36 & 0.76/0.55 & 3.9 \\
\botrule
\end{tabular}
\footnotetext{Batches of 50 consecutive parts, 95\,\% quantile of the difference of their 10\,\% trimmed means; interval from 2000 batch resamples. Ratios per geometry (concave/convex), classical statistics pooled over the series as root mean squares: $\sigma_\mathrm{LT}$, standard deviation of the parts of a series; $\sigma_\mathrm{b}$, between batch means, corrected for $\sigma_\mathrm{w}^2/50$; $\sigma_\mathrm{w}$, within batches. Drift: floor over that of randomly permuted series.}
\end{table}
%% END GENERATED floors

\subsection{What production resolves and how far the simulation misses it}\label{sec:res-sim}

All 63 geometry contrasts exceed the floor and blank-holder force is resolved in 105 of 126 contrasts, but lubrication only in 51 of 126\DIFdelbegin \DIFdel{(42 with a bootstrap probability of at least 0.95)}\DIFdelend , so a lubrication variant is often a near-replicate of its siblings. The minimum resolvable force change is 21--67\,kN between 100 and 300\,kN for most characteristics and lies beyond either force interval for the wall angle and the dome. The nominal simulation misses production by far more than a floor: its median offset reaches 109 floors, its force effect on the draw-in is 2.7 to 5.1 times the measured one, and the cut arms of the concave cups bend downwards in 89--100\,\% of the simulations but upwards in every part. Part of the offset is definitional: the cup depth differs by a median 0.90 and 0.95\,mm, consistent with the scanned surface against the shell mid-surface, and removing this offset alone brings the decisive distance of the nominal simulation from 108 to 65 floors.

\subsection{Decisions by evidence relation}\label{sec:res-decisions}

Table~\ref{tab:distances} and Fig.~\ref{fig:decisions} score the rules. The nominal simulation decides at 108 floors (101--115), erring mostly by rejecting adequate alternatives.

\emph{A new variant with a produced sibling.} Every calibrated rule improves on it, and the simplest is the most decisive: the nearest produced setting, the mean $q_{95}$ of the two siblings, is right beyond 3.4 floors (2.0--4.5), 0.17\,mm of mid-side draw-in; the Gaussian processes decide at 7.2 floors but are safe below one floor.

\emph{A new process setting.} With every alternative at the new design's force held out, all rules lose: the force trend and the nearest produced setting decide at 8.2 and 8.5 floors, the interval rules (M2, M2n, M3, M3n, M5, M5n) only at 13--31, though they are safe at 5.9--16. Interpolating in force is much easier than extrapolating (the nearest setting 4.7 against 9.5 floors), and the scaled simulation, extrapolating a fitted trend, falsely accepts failing alternatives as far as 164 floors from their limit.

\emph{A new design family.} Scored in the floor of the produced family, as the procedure prescribes, every rule without the simulation fails in the two transfers (130--135 floors)\DIFdelbegin \DIFdel{, because }\DIFdelend \DIFaddbegin \DIFadd{: }\DIFaddend nothing in one family's production record describes the other. With the simulation, the bias correction decides at 33 floors and M5 at 35, and M3 is safe at 17 floors, M2 at 26 and M5 at 31; in the floor of the new family, which a team would not have, M2 is safe at 16, M3 at 21 and M5 at 28. The best rule was thus safe at about 16--17 floors in either floor, some 0.8\,mm of mid-side draw-in\DIFdelbegin \DIFdel{or 1.3$^\circ$ of wall angle}\DIFdelend , but two transfers do not resolve which rule is safest: the safe distance of M3 minus that of M2 lies between $-16$ and $+7$ floors in the produced family's floor and between $-3$ and $+8$ in the new family's.

%% BEGIN GENERATED distances
\begin{table}[t]
\tabsize
\setlength{\tabcolsep}{1.9pt}
\caption{Decisive and safe distances (floors) by the relation of the new design to the produced evidence}\label{tab:distances}
\DIFdelbeginFL %DIFDELCMD < \begin{tabular}{@{}lrrrrrrrrrr@{}}
%DIFDELCMD < \toprule
%DIFDELCMD <  & \multicolumn{2}{c}{New variant} & \multicolumn{4}{c}{New process setting} & \multicolumn{2}{c}{New family} & \multicolumn{2}{c}{Pooled} \\
%DIFDELCMD < \cmidrule(lr){2-3}\cmidrule(lr){4-7}\cmidrule(lr){8-9}\cmidrule(lr){10-11}
%DIFDELCMD < Rule & $\delta_\mathrm{d}/\delta_\mathrm{s}$ & FA & $\delta_\mathrm{d}/\delta_\mathrm{s}$ & FA & interp. & extrap. & $\delta_\mathrm{d}/\delta_\mathrm{s}$ & FA & variant & setting \\
%DIFDELCMD < \midrule
%DIFDELCMD < M0 & 108 & 17 & 108 & 17 & 109 & 106 & 104 & 28 & 108 & 108 \\
%DIFDELCMD < M1 & 15 & 15 & 24 & 24 & 11 & 25 & 33 & 35 & 27 & 28 \\
%DIFDELCMD < M1s & 4.9 & 4.7 & 19 & 164 & 19 & 19 & 121 & 125 & 15 & 20 \\
%DIFDELCMD < \quad interval & 4.7--6.3 &  & 14--21 &  & 19--20 & 15--19 & 119--121 &  &  & \\
%DIFDELCMD < M2 & 23/0.15 & 0.30 & 18/8.1 & 8.3 & 18/2.7 & 18/9.9 & 41/26 & 27 & 53/0 & 53/1.0 \\
%DIFDELCMD < \quad interval & 22--23 &  & 17--18 &  & 15--18 & 17--19 & 39--42/9.8--32 &  &  & \\
%DIFDELCMD < M5 & 7.2/0.75 & 0.15 & 19/6.8 & 7.6 & 19/2.8 & 19/8.9 & 35/31 & 33 & 11/0.05 & 28/5.4 \\
%DIFDELCMD < \quad interval & 6.6--8.4 &  & 18--20 &  & 18--20 & 18--20 & 31--39/28--32 &  &  & \\
%DIFDELCMD < M5n & 7.2/0.55 & 0.10 & 17/6.1 & 9.0 & 16/0 & 18/8.6 & 133/128 & 127 & 9.8/0.20 & 16/2.7 \\
%DIFDELCMD < \quad interval & 6.4--7.5 &  & 16--18 &  & 15--16 & 16--19 & 131--133/127--128 &  &  & \\
%DIFDELCMD < NN & 3.4 & 2.8 & 8.5 & 10 & 4.7 & 9.5 & 130 & 129 & 3.4 & 8.5 \\
%DIFDELCMD < \quad interval & 2.0--4.5 &  & 5.5--10 &  & 4.3--5.1 & 6.6--11 & 129--131 &  &  & \\
%DIFDELCMD < \midrule M0w & 119/94 & 9.2 & 119/94 & 9.2 & 119/94 & 119/77 & 117/79 & 10 & 119/94 & 119/94 \\
%DIFDELCMD < Mc & 27 & 28 & 39 & 39 & 14 & 40 & 43 & 53 & 32 & 33 \\
%DIFDELCMD < M3 & 23/0 & 0 & 31/16 & 13 & 29/13 & 33/17 & 40/17 & 18 & 25/0 & 32/5.1 \\
%DIFDELCMD < \quad interval & 17--26 &  & 25--36 &  & 24--31 & 25--40 & 36--45/15--19 &  &  & \\
%DIFDELCMD < M3n & 11/0.05 & 0 & 13/7.4 & 9.2 & 15/7.0 & 13/7.8 & 135/126 & 124 & 15/0 & 15/4.9 \\
%DIFDELCMD < M1n & 9.1 & 10 & 10 & 12 & 4.8 & 12 & 131 & 129 & 130 & 130 \\
%DIFDELCMD < M1sn & 3.5 & 3.5 & 8.2 & 9.0 & 4.7 & 8.9 & 131 & 129 & 77 & 111 \\
%DIFDELCMD < M2n & 19/0 & 0.20 & 15/5.9 & 9.2 & 13/0 & 16/9.1 & 133/128 & 126 & 325/0 & 325/4.0 \\
%DIFDELCMD < M6 & 17/0.05 & 0 & 56/6.2 & 5.8 & 47/0 & 56/8.1 & 35/28 & 28 & -- & -- \\
%DIFDELCMD < \botrule
%DIFDELCMD < \end{tabular}
%DIFDELCMD < %%%
\DIFdelendFL \DIFaddbeginFL \begin{tabular}{@{}lrrrrrrrrrr@{}}
\toprule
 & \multicolumn{2}{c}{New variant} & \multicolumn{4}{c}{New process setting} & \multicolumn{2}{c}{New family} & \multicolumn{2}{c}{Pooled} \\
\cmidrule(lr){2-3}\cmidrule(lr){4-7}\cmidrule(lr){8-9}\cmidrule(lr){10-11}
Rule & $\delta_\mathrm{d}/\delta_\mathrm{s}$ & FA & $\delta_\mathrm{d}/\delta_\mathrm{s}$ & FA & interp. & extrap. & $\delta_\mathrm{d}/\delta_\mathrm{s}$ & FA & variant & setting \\
\midrule
M0 & 108 & 17 & 108 & 17 & 109 & 106 & 104 & 28 & 108 & 108 \\
M1 & 15 & 15 & 24 & 24 & 11 & 25 & 33 & 35 & 27 & 28 \\
\quad interval & 12--18 &  & 21--25 &  & 10--11 & 24--25 & 31--36 &  &  & \\
M1s & 4.9 & 4.7 & 19 & 164 & 19 & 19 & 121 & 125 & 15 & 20 \\
\quad interval & 4.7--6.3 &  & 14--21 &  & 19--20 & 15--19 & 119--121 &  &  & \\
M2 & 23/0.15 & 0.30 & 18/8.1 & 8.3 & 18/2.7 & 18/9.9 & 41/26 & 27 & 53/0 & 53/1.0 \\
\quad interval & 22--23 &  & 17--18 &  & 15--18 & 17--19 & 39--42/9.8--32 &  &  & \\
M5 & 7.2/0.75 & 0.15 & 19/6.8 & 7.6 & 19/2.8 & 19/8.9 & 35/31 & 33 & 11/0.05 & 28/5.4 \\
\quad interval & 6.6--8.4 &  & 18--20 &  & 18--20 & 18--20 & 31--39/28--32 &  &  & \\
M5n & 7.2/0.55 & 0.10 & 17/6.1 & 9.0 & 16/0 & 18/8.6 & 133/128 & 127 & 9.8/0.20 & 16/2.7 \\
\quad interval & 6.4--7.5 &  & 16--18 &  & 15--16 & 16--19 & 131--133/127--128 &  &  & \\
NN & 3.4 & 2.8 & 8.5 & 10 & 4.7 & 9.5 & 130 & 129 & 3.4 & 8.5 \\
\quad interval & 2.0--4.5 &  & 5.5--10 &  & 4.3--5.1 & 6.6--11 & 129--131 &  &  & \\
\midrule M0w & 119/94 & 9.2 & 119/94 & 9.2 & 119/94 & 119/77 & 117/79 & 10 & 119/94 & 119/94 \\
Mc & 27 & 28 & 39 & 39 & 14 & 40 & 43 & 53 & 32 & 33 \\
M3 & 23/0 & 0 & 31/16 & 13 & 29/13 & 33/17 & 40/17 & 18 & 25/0 & 32/5.1 \\
\quad interval & 17--26 &  & 25--36 &  & 24--31 & 25--40 & 36--45/15--19 &  &  & \\
M3n & 11/0.05 & 0 & 13/7.4 & 9.2 & 15/7.0 & 13/7.8 & 135/126 & 124 & 15/0 & 15/4.9 \\
M1n & 9.1 & 10 & 10 & 12 & 4.8 & 12 & 131 & 129 & 130 & 130 \\
M1sn & 3.5 & 3.5 & 8.2 & 9.0 & 4.7 & 8.9 & 131 & 129 & 77 & 111 \\
\quad interval & 3.1--4.2 &  & 6.7--9.0 &  & 4.3--5.1 & 8.6--9.5 & 129--131 &  &  & \\
M2n & 19/0 & 0.20 & 15/5.9 & 9.2 & 13/0 & 16/9.1 & 133/128 & 126 & 325/0 & 325/4.0 \\
M6 & 17/0.05 & 0 & 56/6.2 & 5.8 & 47/0 & 56/8.1 & 35/28 & 28 & -- & -- \\
\botrule
\end{tabular}
\DIFaddendFL \footnotetext{Over all seven characteristics; one number for rules that always decide. FA: safe distance on the false-accept side (failing alternatives). Interval: bootstrap 95\,\% interval of the decisive distance and, for a new family, of the safe distance (fits fixed). 126 cases per relation (interp.: 42 interpolated; extrap.: 84 extrapolated); a new family rests on two transfers and is scored in the floor of the produced family (Section~\ref{sec:floor}). Pooled: the other geometry's alternatives added to the calibration set. Below the line, the further rules (M1sn: the force trend, the counterpart of M1s without the simulation) and, as a robustness variant, M6 (the M5 mean with a conformal margin on standardised leave-one-out residuals).}
\end{table}
%% END GENERATED distances

\begin{figure}[t]
\centering
\includegraphics[width=\textwidth]{F3_decisions_compare}
\DIFdelbeginFL %DIFDELCMD < \caption{Verdict shares against the requirement distance $d$ (symmetric logarithmic axis, linear within $\pm$10 floors; 126 cases per panel; a new family in the produced family's floor), with the decisive and safe distances}%%%
\DIFdelendFL \DIFaddbeginFL \caption{Verdict shares against the requirement distance $d$ (symmetric logarithmic axis, linear within $\pm$10 floors; 126 cases per panel; a new family in the produced family's floor), with $\delta_\mathrm{d}$ and $\delta_\mathrm{s}$}\DIFaddendFL \label{fig:decisions}
\end{figure}

\subsection{Where the distances come from}\label{sec:res-drivers}

Three features of the design shape these distances. First, the sibling relation mixes replication with variation: in 58 of the 126 cases neither sibling differs resolvably from the held-out alternative, and the nearest setting is then right beyond 0.85 floors, the reproducibility of $q_{95}$ between near-replicate runs; it needs 2.5 floors when one sibling differs resolvably (34 cases) and 4.9 when both do (34 cases), where the force trend needs 4.5 and M5 9.7. Second, the extrapolation is confounded with stroke speed: the nearest setting decides at 11 floors at 100\,kN, where the press also ran slower, against 4.9 at 500\,kN (the force trend at 8.9 at both). Taken together, the nearest setting decided at \DIFaddbegin \DIFadd{up to }\DIFaddend about 5 floors whenever the new design differed resolvably and force and speed did not change together, though each stratum has 34 to 58 cases and each distance rests on a handful of verdicts. Third, the other family's record does not help with a produced sibling: pooling it makes most rules less decisive (M5 from 7.2 to 11 floors), though it raises M5's coverage from 86 to 90\,\%.

The relation therefore does not govern decision fitness in general. Over the eleven calibrated rules and three relations, the relation explains 67\,\% of the variation of the log decisive distance and the rule 6\,\% (69 and 6\,\% in the new family's own floor), but only because rules without the simulation fail for a new family by construction; over the two relations within a family, the rule explains 68\,\% and the relation 15\,\% (71 and 15\,\% with M1sn). Pooled over strata, the best attainable distance grows with novelty: 3.4 floors with a produced sibling, 8.2 for a new setting, 33 for a new family.

\subsection{What the simulation and learning contribute}\label{sec:res-ml}

Pairing each calibrated rule with its counterpart without the simulation isolates what the simulation adds as used. Added as an offset, it adds nothing with a produced sibling and degrades two rules: M5 decides where M5n does (difference 0.0 floors, $-0.4$ to $+1.1$), while M2 and M1 decide 3.9 ($+0.5$ to $+7.4$) and 6.2 floors ($+4.6$ to $+9.7$) further out than M2n and M1n, and M1 13.5 further out for a new setting, because the simulated force effect is too large. Rescaled, it does no better: within a geometry the nominal simulation is a function of the force, and the straight line in force (M1sn) decides 1.4 floors closer than the scaled simulation with a sibling ($+0.65$ to $+3.0$) and 11 closer for a new setting ($+5.3$ to $+13.8$), in every resample and refit configuration; the 4.2 floors that the scaled simulation gains over the median of the produced alternatives with a sibling come from the force level. Only for a new family is the simulation decisive, by 92 to 98 floors as an offset and 10 rescaled, and even there production knowledge carries over where a drawing nominal transfers: for the wall angle, the design angle plus the produced family's deviation from its own decides at 1.9 floors, against 7.5 for M5 (for the arm angle and the dome, whose nominal is zero in both families, this baseline is M1n, at 44 and 27 floors). A Gaussian-process surrogate of the simulations over force, friction and thickness, with a leave-one-out root-mean-square error below 2 floors for 12 of the 14 characteristics and geometries, decides as the simulation does (109 against 108 floors as M0, 15 against 15 as M1 with a sibling).

Learning adds intervals rather than accuracy. Used as point rules, the centres of the M5, M5n and M3n intervals decide at 3.3, 3.1 and 3.7 floors with a produced sibling, as close as the nearest setting (Table~\ref{tab:intervals}). With a sibling the Gaussian-process intervals cover 85--86\,\% of the true $q_{95}$ against a nominal 90\,\%, but their noise variance sits at its lower bound, the variance between the calibration series of one setting, in 97--99\,\% of these fits and in every new-setting and new-family fit, so this coverage comes from the quasi-replicate series rather than from the learner; the largest-residual margin covers 89\,\% and the jackknife+ intervals 96\,\% (M3) and 89\,\% (M3n), above their guarantee of 7/9. For a new setting the calibrated intervals cover 43--55\,\% (M2n and M5n 93\,\% interpolating), for a new family the Gaussian processes at most 12\,\%: trees do not extrapolate, two force levels identify at most a linear trend, and in the median new-setting fit the length scale of the force sits at its lower bound.

\subsection{Calibration budget}\label{sec:res-budget}

A single produced sibling decides (Fig.~\ref{fig:budget}): with one alternative at the new design's force, the bias correction, the envelope rule and the nearest setting coincide and decide at 3.5 floors (2.3--4.6) without a wrong verdict at 10 floors. Further alternatives do not improve the nearest setting; they make the bias correction worse (12--18 floors), the envelope rule safer but less decisive, and M5, shown from three alternatives because a Gaussian process fits its hyperparameters to them, safer and more decisive (13 to 7.2 floors). Without a sibling, bracketing the new force buys safety but not decisiveness: the envelope rule and M5 are safe at 2.7--4.8 floors but decide only at 15--25, against 4.7--6.0 for the nearest setting, and extrapolation stays unsafe at every budget. \DIFdelbegin \DIFdel{Averaged over all subsets, the distances shrink with the budget mainly because the share of subsets containing a sibling grows, from 46\,\% with two alternatives to 96\,\% with six. }\DIFdelend With calibration series of 50, 100 or 250 parts instead of 500, no distance within a family moves by more than 1.5 floors, and for a new family by at most 4 \DIFaddbegin \DIFadd{in its own floor }\DIFaddend (Table~\ref{tab:robust}): rules can be calibrated on short runs, \DIFdelbegin \DIFdel{here scored against the truths and floors of }\DIFdelend \DIFaddbegin \DIFadd{but the floors and the held-out truths came from }\DIFaddend the full series, and \DIFdelbegin \DIFdel{only the floor needs a run long enough to show the drift}\DIFdelend \DIFaddbegin \DIFadd{whether truths from short runs suffice to qualify the rules was not tested}\DIFaddend .

\begin{figure}[t]
\centering
\includegraphics[width=\textwidth]{F4_budget}
\caption{Decisive and safe distances against the number of produced alternatives of the family, by the relation of the new design to them; bands: 95\,\% intervals over the held-out alternatives}\label{fig:budget}
\end{figure}

\subsection{The procedure's decision rule}\label{sec:res-procedure}

What a designer observes for a single verdict is the margin between the predicted interval and the requirement, and its value as evidence depends on the relation (Fig.~\ref{fig:reliability}): with a produced sibling at least 95\,\% of the decided verdicts of the interval rules are correct at every margin, for a new setting the share grows with the margin, and for a new family M2 and M5 are right in at most 79\,\% of their decided verdicts.

Step 5 turns this into guard bands (Table~\ref{tab:step6}). Under the reference prior the interval rules need none with a sibling and the nearest setting 0.25 floors; for an interpolated setting the nearest setting needs 2.0 floors and then decides at 6.7 floors without trials at 10 floors, whereas M2 and M5 need 0.75 and 0.5 floors but send 31 and 40\,\% to trial; for an extrapolated setting it needs 7.0 floors and sends a quarter to trial. M0 and M0w never qualify, nor do M1 and Mc with a sibling, M1s for an interpolated setting, M1, M1s, M2, M5 and Mc for an extrapolated one, or any rule for a new family.

These bands belong to the prior as much as to the rules. A prior concentrated near the truth, where release-level requirements lie (the grid within $\pm5$ floors), raises the band of the nearest setting to 2.0 floors with a sibling, 4.25 interpolating and 14.25 extrapolating \DIFaddbegin \DIFadd{(8.0, 5.0 and 14.75 within $\pm3$)}\DIFaddend ; a uniform prior within $\pm20$ floors lowers them to 0, 0.25 and 3.5 floors and qualifies M2 and M5 for an extrapolated setting (6.0 and 5.0 floors), and one within $\pm50$ floors qualifies M2, M3 and M5 for a new family (6.25, 6.5 and 15 floors). That every new-family design goes to trial thus follows from the reference prior and the 20-floor cap, as the practice of trying out every new die would have it. The bands are chosen on the evaluation data, which makes the trial rates optimistic\DIFdelbegin \DIFdel{, and the guards}\DIFdelend \DIFaddbegin \DIFadd{. The guards, reported as diagnostics, }\DIFaddend refuse indiscriminately: \DIFaddbegin \DIFadd{the range guard refuses exactly the extrapolated designs (84 of the 126 new-setting cases), and at requirements 10 floors from the truth }\DIFaddend 96\,\% of \DIFdelbegin \DIFdel{the range guard's refusals for a new setting }\DIFdelend \DIFaddbegin \DIFadd{its refusals }\DIFaddend would have been right for the nearest setting and 72\,\% for M5 (novelty guard 96 and 73\,\%).

\emph{A worked example.} A tryout engineer has to decide whether the concave cup at 300\,kN with medium lubrication, a new setting, can be released without a trial run, with the six alternatives at 100 and 500\,kN produced and a requirement, chosen for illustration, that 95\,\% of the parts have a mid-side draw-in, by which tryout sets the blank-holder force, of at most 15.30\,mm. The floor is 0.050\,mm, the relation is interpolation and no guard fires. Alone, the nominal simulation predicts 15.66\,mm and fails the design. The nearest produced setting, the mean of the six alternatives at the neighbouring forces, predicts 14.99\,mm, a margin of 6.2 floors, more than its guard band of 2.0 floors (5.0 under the grid within $\pm3$ floors), so step 6 returns meets. Widened by their guard bands, M2, M5 and M5n return 14.41--15.41, 14.23--15.61 and 14.66--15.33\,mm and send the design to trial. The 500 parts had a 95th percentile of 14.88\,mm, 8.4 floors below the limit. The case is typical, since at 8.5 floors above the truth step 6 with the nearest setting is right in all 42 interpolated cases, but its guard band was estimated on data that include it.

\begin{figure}[t]
\centering
\includegraphics[width=\textwidth]{F5_reliability}
\DIFdelbeginFL %DIFDELCMD < \caption{Share of correct verdicts among the decided ones with at least the given margin between interval and requirement, under the reference prior (47 grid distances within $\pm$20 floors, equal weights); filled markers: guard bands}%%%
\DIFdelendFL \DIFaddbeginFL \caption{Share correct among decided verdicts with at least the given margin; reference prior: 47 grid distances within $\pm$20 floors, equal weights; new family in the produced family's floor; filled markers: $g$}\DIFaddendFL \label{fig:reliability}
\end{figure}

\subsection{Requirements anchored in capability and in tolerances}\label{sec:res-requirements}

Upper limits at which an alternative has a performance index of 1.0, 1.33, 1.67 or 2.0 lie a median 1.0, 1.7, 2.5 and 3.1 floors above its 95th percentile; up to an index of 1.33 they lie inside every decisive distance except a near-replicate sibling's (0.85 floors). The test is one-sided by construction: every alternative meets these limits, so a verdict can only be right, a false reject or a trial, and a rule biased towards meets scores well. At an index of 1.33 the shares of right verdicts, false rejects and trials are, with a produced sibling and for a new setting, 91/9/0 and 65/35/0\,\% for the nearest setting, 44/2/53 and 32/16/52\,\% for M5, and 81/19/0 and 79/21/0\,\% for the scaled simulation, which falsely accepts failing alternatives elsewhere; for a new family no rule is right in more than 54\,\%. Release-level requirements are thus decided at the design stage only with a produced sibling, and then by the production record. General angular tolerances of ISO~2768-1 \citep{iso1989tolerances} give a coarser scenario: only 32 of its 108 decisions lie within 5 floors of the limit, 6 of them on failing alternatives, and a cost comparison inherits this coarseness. With a false reject at half the cost of a false accept, the cheapest rule is (i) with a produced sibling, the nearest setting, the scaled simulation and the force trend, tied at zero cost; (ii) for a new setting, the force trend at every trial cost\DIFaddbegin \DIFadd{, by three false rejects of the nearest setting among 108 decisions}\DIFaddend ; and (iii) for a new family, a trial of every design while a trial costs at most a tenth of a false accept, the envelope rule at a fifth and M5 beyond.

\subsection{Robustness}\label{sec:res-robust}

Table~\ref{tab:robust} repeats the analysis under variants of the data, constants and model specifications. Adjusting for punch temperature changes no calibrated distance by more than 2.6 floors. Dropping the warm-up of the first 150 parts shrinks the floors and so lengthens the distances, and floors taken between series shorten them. Other floor protocols rescale the decisive distances by 0.60 to 1.42, and the Gaussian-process specifications move them by at most 1.2 floors and the safe distances by at most 2.1. In every variant that includes them, the most decisive rule within a family is the nearest setting or the force trend, and with a produced sibling the nearest setting (tied once). The refitting bootstrap widens the intervals of the learned rules most (M5 4.0--11 floors with a sibling, the nearest setting 2.3--4.6), yet in all 200 resamples (48 configurations) the nearest setting is more decisive than M5 and M5n, with a sibling and for a new setting alike, and the force trend than the scaled simulation (by 1.35--3.6 and 8.8--12 floors); point estimates at the edge of their refit intervals (M1s with a sibling at 4.9 floors, 4.9--7.3, and the safe distances of M2 and M5) reflect what losing a pattern does.

%% ----------------------------------------------------------------------
\section{Discussion}\label{sec:discussion}

\subsection{What the evaluation establishes}

Five findings answer the research questions. Two follow from the definitions and hold wherever the framework is applied: the floor puts every rule and characteristic on a common production-referenced scale, though not on a common scale of nonconforming fractions, and as a within-run quantity it bounds the variation of series production from below (RQ1); and an interval rule trades decisiveness for safety, which the two distances make explicit, while its distribution-free intervals are guaranteed their level only where the calibration alternatives are exchangeable with the new design. Three are observations on this dataset and hypotheses for confirmation: the best attainable distance grows with the novelty of the design, from 3.4 floors with a produced sibling (the nearest setting 0.85 for near-replicates, 4.9 for resolvable siblings) to 8.2 for a new setting (\DIFaddbegin \DIFadd{the force trend; by stratum the nearest setting needed }\DIFaddend 4.7 interpolated \DIFdelbegin \DIFdel{; 4.9 and 8.9 }\DIFdelend \DIFaddbegin \DIFadd{and 4.9 }\DIFaddend extrapolated to 500\DIFdelbegin \DIFdel{and }\DIFdelend \DIFaddbegin \DIFadd{\,kN, the force trend 8.9 to }\DIFaddend 100\,kN\DIFaddbegin \DIFadd{, where the nearest setting needed 11}\DIFaddend ) and 33 for a new family, while within a family the choice of rule matters more than the relation (RQ2); production evidence buys safety before decisiveness, and interval rules held their level with a produced sibling but not when extrapolating or for a new family; and with this simulator and six to nine produced alternatives at two or three force levels, the simulation, as an offset or rescaled, contributed only across families, and \DIFdelbegin \DIFdel{learned models }\DIFdelend \DIFaddbegin \DIFadd{M5, M5n and M3n }\DIFaddend predicted as well as the nearest produced setting with a produced sibling while adding abstention (RQ3).

\begin{table}[t]
\tabsize
\setlength{\tabcolsep}{4pt}
\caption{Guidance for design-stage manufacturability decisions, with its basis}\label{tab:guidelines}
\DIFdelbeginFL %DIFDELCMD < \begin{tabular}{@{}L{2.9cm}L{7.75cm}L{1.85cm}@{}}
%DIFDELCMD < \toprule
%DIFDELCMD < Situation & Guidance & Basis \\
%DIFDELCMD < \midrule
%DIFDELCMD < \multicolumn{3}{@{}l}{\emph{From the framework (prescriptive; usefulness not evaluated)}} \\
%DIFDELCMD < Before relying on a predictor & Measure the floor on consecutive production of the family with a stated protocol and report $\sigma_\mathrm{b}$, $\sigma_\mathrm{w}$ and its lag dependence; treat it as a lower bound of series variation. & framework \\
%DIFDELCMD < Qualifying a rule & Evaluate it for the evidence relation of the coming designs by grouped cross-validation, and split near-replicate siblings from resolvable ones. & framework \\
%DIFDELCMD < Choosing a rule & Compare it with the nearest produced setting and with its own version without the simulation, on decisive and safe distance jointly or by expected cost. & framework \\
%DIFDELCMD < Deciding a single design & Derive the rule's guard band for the distribution of requirement distances the team expects and apply it to the observable margin; a margin inside it, or a fired guard, means a trial. & framework \\
%DIFDELCMD < Planning a requirement & As a heuristic, plan a trial if the expected margin between requirement and predicted 95th percentile is below the safe distance of the best qualified rule; no decision is reliable closer than a trial reproduces the quantile. & framework \\
%DIFDELCMD < \multicolumn{3}{@{}l}{\emph{Observed on this dataset (18 alternatives, 2 geometries, 7 characteristics)}} \\
%DIFDELCMD < New variant with a produced sibling & nearest produced setting right beyond 3.4 floors (0.85 for near-replicates, 4.9 for resolvable siblings); interval rules safe below one floor, coverage 85--96\,\% & 126 cases \\
%DIFDELCMD < New process setting & nearest setting or force trend right beyond 4.7 floors interpolated, 4.9 extrapolated to 500\,kN and 8.9 to 100\,kN (with a speed change); interval rules safe at 5.9--16 floors, coverage 43--55\,\% & 42 + 84 cases \\
%DIFDELCMD < New family & best rule safe at 16--17 floors in either family's floor, which rule not resolved; rules without the simulation fail; a transferable drawing nominal helps & two transfers \\
%DIFDELCMD < Release-level requirements & lie 1.7 floors above the 95th percentile at an index of 1.33, inside every decisive distance except a near-replicate's; decided only with a produced sibling & 126 cases \\
%DIFDELCMD < \botrule
%DIFDELCMD < \end{tabular}
%DIFDELCMD < %%%
\DIFdelendFL \DIFaddbeginFL \begin{tabular}{@{}L{2.9cm}L{7.75cm}L{1.85cm}@{}}
\toprule
Situation & Guidance & Basis \\
\midrule
\multicolumn{3}{@{}l}{\emph{From the framework (prescriptive; usefulness not evaluated)}} \\
Before relying on a predictor & Measure the floor on consecutive production of the family with a stated protocol and report $\sigma_\mathrm{b}$, $\sigma_\mathrm{w}$ and its lag dependence; treat it as a lower bound of series variation. & framework \\
Qualifying a rule & Evaluate it for the evidence relation of the coming designs by grouped cross-validation, and split near-replicate siblings from resolvable ones. & framework \\
Choosing a rule & Compare it with the nearest produced setting and with its own version without the simulation, on decisive and safe distance jointly or by expected cost. & framework \\
Deciding a single design & Derive the rule's guard band for the distribution of requirement distances the team expects and apply it to the observable margin, in place of range and family guards; a margin inside it means a trial. & framework \\
Planning a requirement & As a heuristic, plan a trial if the expected margin between requirement and predicted 95th percentile is below the safe distance of the best qualified rule; no decision is reliable closer than a trial reproduces the quantile. & framework \\
\multicolumn{3}{@{}l}{\emph{Observed on this dataset (18 alternatives, 2 geometries, 7 characteristics)}} \\
New variant with a produced sibling & nearest produced setting right beyond 3.4 floors (0.85 for near-replicates, 4.9 for resolvable siblings); interval rules safe below one floor, coverage 85--96\,\% & 126 cases \\
New process setting & nearest setting right beyond 4.7 floors interpolated and 4.9 extrapolated to 500\,kN, force trend 8.9 to 100\,kN (nearest setting 11; speed change); interval rules safe at 5.9--16 floors, coverage 43--55\,\% & 42 + 84 cases \\
New family & best rule safe at 16--17 floors in either family's floor, which rule not resolved; rules without the simulation fail; a transferable drawing nominal helps & two transfers \\
Release-level requirements & lie 1.7 floors above the 95th percentile at an index of 1.33, inside every decisive distance except a near-replicate's; decided only with a produced sibling & 126 cases \\
\botrule
\end{tabular}
\DIFaddendFL \end{table}

\subsection{Using the framework in design}

As a planning heuristic (Table~\ref{tab:guidelines}), the expected margin between a requirement and the predicted 95th percentile can be compared with the safe distance of the best rule qualified for the relation; if it is smaller, a trial or more production evidence has to be planned, and the decision itself needs the guard band, because the safe distance is defined on the true margin. Release-level requirements make this the usual case, as the practice of a production trial run for a process change expects: the rules decide reliably where a trial is cheapest, a new force on an existing tool being a short press run, and none does near release-level limits where a trial is dearest, for a new family that may need a new tool. The most decisive rule is arguably also the most transparent: the nearest produced setting points to the produced alternatives whose record it averages. The variants that steps 1 to 5 need come from process development and tryout, whose short series of several settings support the sibling and new-setting relations, and from a family's production history across part numbers, which supports the new-family relation; production part approval, whose initial process study covers the released setting, supplies the floor but not the variants.

\subsection{Simulation and machine learning in this regime}

Within a design family the production record carried the decision, even against a rescaled simulation. The simulated force trend of the two draw-ins was 2.7 to 5.1 times too large, and calibrating the friction coefficient did not repair it, which points to the friction law and the blank-holder model rather than to a parameter. A multi-task model cannot help a family with no produced alternative, because the correlation between families cannot be estimated without its data; in the pooled scopes the Gaussian process with a 0/1 geometry input is a restricted intrinsic coregionalisation model (equal variances, one positive correlation between the families, shared length scales), which improved coverage but not decisiveness. These findings concern one simulator of low fidelity, used through offsets and one scale, and learners with six to nine produced alternatives at two or three force levels. For evaluating AI-based DFM support the lesson is methodological: a learned model should be compared with a production-record baseline, scored by evidence relation in a unit that production defines, and shown to its user with a guard band. A study with design teams could test two hypotheses on reliance \citep{lee2004trust}, with cases of known production outcome: designers shown the guard-banded verdict accept a smaller share of wrong verdicts within the safe distance than designers shown the verdict alone, at an equal rate of right acceptances; and, as the same AI assistance helped some teams and hindered others \citep{zhang2021cautionary}, the guard band narrows the spread of that share between teams.

\subsection{Implications for design research}

\DIFdelbegin \DIFdel{The results give constructs of design research a production-referenced measure. }\DIFdelend For variant and platform design \citep{pahl2007engineering, simpson2006product} and the capability databases that serve them \citep{tata1999process, thornton2004variation}, the results argue for indexing production records by their relation to the coming design rather than for a more sophisticated model on top. The safe distance is the part of a margin that covers predictive uncertainty, as distinct from the part that absorbs change \citep{eckert2017safety, eckert2019design}. The growth of the decisive distance\DIFaddbegin \DIFadd{, pooled over strata, }\DIFaddend from about 3 to 8 and over 30 floors says for which overlaps of development \citep{krishnan1997model, terwiesch2002exchanging} preliminary manufacturability information can be released, the new-family safe distance, some 0.8\,mm of draw-in, how far a set-based design \citep{sobek1999toyota} can narrow a feasible set before production evidence exists, and the distances and trial rates when a test is cheaper than a decision \citep{thomke2001sequential, salado2018mathematical}. In the terms of the design research methodology \citep{blessing2009drm}, the evaluation shows structural validity, that the metrics discriminate between rules and relations on real data \citep{seepersad2006validation}; performance validity would require better decisions with the framework.

\begin{table}[t]
\tabsize
\setlength{\tabcolsep}{3pt}
\caption{What the framework requires, mapped onto other processes (an untested proposal)}\label{tab:processes}
\begin{tabular}{@{}L{1.9cm}L{2.4cm}L{2.5cm}L{2.7cm}L{2.7cm}@{}}
\toprule
Process & Batch and run & Main variation & Evidence relations & Floor \\
\midrule
Sheet forming (this paper) & consecutive parts of one series & drift in a run; coil, die state & variant, setting, family & batch centres in a run \\
Injection moulding & shots of one cavity and set-up & cavity, shot, set-up, material lot & cavity, process window, new mould & per cavity; cavity offsets as effects \\
Machining & consecutive parts between offsets & tool wear, offset-compensated & tool path, cutting data, new part & batch centres within an offset interval \\
Casting & castings of one pattern and melt & melt, pattern position & alloy, gating, new pattern & within a melt; melts as runs \\
Additive manufacturing & replicates in one build & build, position, orientation & orientation, settings, new part & within a build, via replicated anchors \\
\botrule
\end{tabular}
\end{table}

\subsection{Reuse for other processes and AI-based DFM tools}\label{sec:reuse}

The framework needs several produced variants of a design family whose relations span those of the coming designs, a series of consecutive parts or replicated batches of one design per variant, a design-stage predictor for every variant, and characteristics measured identically on parts and predictions; here rules calibrated on about 50 parts per variant scored \DIFdelbegin \DIFdel{as well as with }\DIFdelend \DIFaddbegin \DIFadd{within 1.5 floors of those calibrated on }\DIFaddend full series, \DIFdelbegin \DIFdel{and only the floor needed long runs}\DIFdelend \DIFaddbegin \DIFadd{against truths and floors from the full series}\DIFaddend . Attribute outcomes such as splits need a requirement on a rate, which the framework does not yet cover. Other predictors need three adaptations. A classifier with a fixed threshold does not take the requirement as input, so its operating characteristic is estimated by binning test cases by their true margin, not by the exact estimator of Section~\ref{sec:distances}, unless the requirement is made an input; a probabilistic classifier needs two thresholds to return a trial. A language-model check may return verdicts that are not monotone in the requirement and that vary between queries, so $\kappa$ and $\omega$ are computed on a requirement grid with repeated queries, and the share of non-monotone verdict sequences is reported. For continuous design spaces the relation becomes a distance from the produced designs in descriptor space or the position relative to their hull, and each stratum needs about as many cases as a relation here: with 126 cases, some 240 verdicts, the 95\,\% threshold already rests on about a dozen. For a generative tool, its manufacturability check is scored. A minimal report comprises
\begin{enumerate}
\item the floor with its protocol, variance components and lag dependence;
\item decisive and safe distances by relation, with refit intervals and one-sided values;
\item coverage and half-width of interval rules, and the source of their noise variance;
\item the nearest-produced-setting and no-simulation baselines;
\item the guard bands with their prior, the trial rates and the performance of the guards;
\item for language-model tools, the prompt and the query protocol.
\end{enumerate}
Table~\ref{tab:processes}, an untested proposal, maps these requirements onto other processes; distances in floors are comparable within a protocol, not across processes, since a floor between builds, cavities or offset intervals is a between-unit quantity (here the floor between series is 1.5 to 9.1 times that within a run). An attempt on open powder bed fusion data of PA12 \citep{leirmo2021data} failed on the second requirement: orientation classes mixed orientations in a batch, eight identical anchors per build were too few for a floor within a build, and three builds give three differences between builds.

\subsection{Threats to validity}

The evaluation covers one material, one press and two geometries, with relations within a family from a $3\times3$ factorial in which lubrication is mostly not resolved; the procedure is general, the numbers are not. The new family rests on two transfers between geometries that differ in three features at once, and the intervals of its distances hold the two fits fixed. Each alternative was produced once, so the floor covers drift within a run but not between days, coils or tool states, the reference $q_{95}$ is that of one run, and the order and dates of the series are unknown; between-series variation bounds the missing component at 1.2--9.8 floors, and scanner drift cannot be separated from process drift. The extrapolation to 100\,kN is confounded with stroke speed, which the simulation omits. The requirements are one-sided limits on the 95th percentile of a characteristic averaged over sides, below the conformance of series release. Several rules, the grouped schemes, the guard bands and their prior were added or chosen after earlier results had been seen and are evaluated on the same data, so the comparison between rules carries selection optimism \citep{cawley2010over}; a single analyst carried out the analysis, with code and data open for replication.

%% ----------------------------------------------------------------------
\section{Conclusions}\label{sec:conclusion}

Design-stage manufacturability decisions can be evaluated against production on a scale production defines: the decisive and safe distances in floors tell how close to a requirement a rule decides and how close it can be relied on not to mislead, and a guard band derived for the requirements a team expects turns them into a rule for single decisions. On 18 deep-drawn alternatives with matched simulations, the best attainable distance grew with novelty: 3.4 floors with a produced sibling (the nearest setting 0.85 for near-replicates, 4.9 for resolvable ones), 8.2 for a new process setting (\DIFaddbegin \DIFadd{the force trend; the nearest setting }\DIFaddend 4.7 interpolated \DIFdelbegin \DIFdel{; }\DIFdelend \DIFaddbegin \DIFadd{and }\DIFaddend 4.9 extrapolated to 500\,kN\DIFdelbegin \DIFdel{and }\DIFdelend \DIFaddbegin \DIFadd{, the force trend }\DIFaddend 8.9 to 100\,kN, with a stroke-speed change), and no rule safe closer than 16--17 floors for a new family, in either family's floor. Within a family the choice of rule mattered more than the relation; \DIFdelbegin \DIFdel{release-level requirements }\DIFdelend \DIFaddbegin \DIFadd{requirements at a performance index of 1.33 }\DIFaddend lay inside every decisive distance except a near-replicate's and need a trial without a produced sibling; and in this small-data regime the simulation as used, offset or rescaled, contributed only across families, while learned 90\,\% intervals covered 85--96\,\% of truths with a produced sibling, under 50\,\% extrapolating or for a new family.

\DIFaddbegin \DIFadd{For design practice, the production record carried the decision within a family, which argues for indexing such records by relation to the coming design rather than for a more elaborate model; across families only rules using the simulation carried the change of geometry, and none decided near release-level limits. For AI-based manufacturability support, the framework offers a standard of evaluation rather than a verdict on machine learning: a tool is scored by evidence relation in floors against a production-record baseline and reported with the checklist of Section~\ref{sec:reuse}.
}

\DIFadd{The evidence is one process and two geometries, and the comparison of rules is exploratory. }\DIFaddend Four extensions follow: series from several days, coils and tool states; families with continuous geometric descriptors; a confirmatory test of the frozen procedure on other data; and a study of whether design teams decide better with it.

%% ----------------------------------------------------------------------

\FloatBarrier
\begin{appendices}
\renewcommand{\floatpagefraction}{0.4}
\renewcommand{\theHtable}{\thetable}% distinct hyperlink targets for Tables A1 and A2

\section{Notation, further results and computation details}\label{app:appendix}

Table~\ref{tab:notation} summarises the notation of Section~\ref{sec:method}, Table~\ref{tab:robust} gives the distances under the robustness variants, Table~\ref{tab:intervals} the coverage, width and centre of the interval rules, and Table~\ref{tab:step6} the procedure's decision rule by relation. Table~\ref{tab:mapping} traces every other number quoted in the text to the script that computes it and the result file that holds it in the archived release \citep{canseven2026archive}.








\emph{Computation details.} Batches need at least 80\,\% valid values. M5 and M5n are Gaussian-process regressions (constant times anisotropic squared-exponential kernel plus white noise, normalised targets, hyperparameters by maximum marginal likelihood without hyperpriors, three restarts, length scales in $[0.1, 100]$) on a geometry indicator, the force in 100\,kN and the lubrication rank; the lower bound of the noise variance, the larger of the mean sampling variance of the calibration alternatives' $q_{95}$ and the pooled variance between their series of one geometry and force, is expressed on the normalised scale. NN takes the mean $q_{95}$ at the nearest force of the family, both neighbours when equally near; M1sn fits $q_{95}=a+b\,$force by least squares. Mc calibrates, per lubrication pattern, the friction coefficient whose simulated draw-ins are closest to the calibration centres. M3 and M3n use histogram gradient boosting (depth 3, 200 iterations, rate 0.08) on geometry, force, lubrication, sheet thickness and oil film, with 500 incoming conditions drawn from calibration parts of the same pattern; their jackknife+ interval takes the leave-one-out predictions for the new alternative, shifted by the median leave-one-out residual, minus and plus the absolute residuals about that median, with its bounds at the extreme order statistics when $n$ is too small for the 0.9 quantiles. Guard bands lie on a grid of 0.25 floors. The refitting bootstrap draws the three lubrication patterns of each geometry with replacement 200 times, at least two of them distinct, and does not restrict the noise floor to distinct series. Software: Python 3.11 with NumPy, pandas, SciPy, scikit-learn \citep{pedregosa2011scikit} and Matplotlib.

\begin{table}[!ht]
\tabsize
\setlength{\tabcolsep}{4pt}
\caption{Notation and terms}\label{tab:notation}
\begin{tabular}{@{}L{2.55cm}L{10.2cm}@{}}
\toprule
Symbol, term & Meaning \\
\midrule
$a$, $c$, $q_p(a,c)$ & alternative (design--process combination), characteristic, $p$-quantile of the parts of $a$ ($p=0.95$) \\
$R$, $d$ & requirement (upper limit on $c$ or $|c|$); its distance from the truth in floors, $R=q_p+dF_c$ \\
$[L_r,U_r]$, $V_r$ & interval of rule $r$ for $q_p$; verdict meets, fails or trial (Eq.~\ref{eq:verdict}) \\
$B$, $m$, $F_c$ & batch size (50); batch centre (10\,\% trimmed mean); production floor (Eq.~\ref{eq:floor}) \\
$\sigma_\mathrm{b}$, $\sigma_\mathrm{w}$ & standard deviation between batch means and between parts within a batch \\
ESR, MRC & effect-to-scatter ratio (Eq.~\ref{eq:esr}); minimum resolvable change, $F_c/|s_c|$ \\
$u$, $v$ & distances in floors by which an interval excludes the truth from above and below \\
$\kappa(\delta)$, $\omega(\delta)$ & shares of correct and wrong verdicts among the feasible verdicts at $|d|=\delta$ \\
$\delta_\mathrm{d}$, $\delta_\mathrm{s}$ & decisive distance ($\kappa\ge0.95$ beyond it) and safe distance ($\omega\le0.05$ beyond it) \\
FA, FR side & one-sided distances for failing alternatives (false accepts) and adequate ones (false rejects) \\
$g$ & guard band: observable margin beyond which at least 95\,\% of the decided verdicts were correct, for a stated prior of requirement distances, here the 47 grid distances within $\pm20$ floors with equal weights (Section~\ref{sec:procedure}); not a property of rule and relation alone \\
Calibration alternatives & produced alternatives whose record calibrates a rule \\
Evidence relation & sibling (same setting produced), new process setting (interpolated, extrapolated), new family \\
Calibration budget & distances against the number $k$ of produced alternatives, by relation \\
Guard & family, range and novelty checks that send a design outside the calibration to trial \\
Measurement resolution & step between adjacent values that the measuring system reports (gauge sense) \\
\botrule
\end{tabular}
\end{table}

%% BEGIN GENERATED robust
\begin{table}[!ht]
\tabsize
\DIFdelbeginFL %DIFDELCMD < \setlength{\tabcolsep}{1.6pt}
%DIFDELCMD < %%%
\DIFdelendFL \DIFaddbeginFL \setlength{\tabcolsep}{1.4pt}
\DIFaddendFL \caption{Decisive/safe distances (floors) under the robustness variants}\label{tab:robust}
\DIFdelbeginFL %DIFDELCMD < \begin{tabular}{@{}lrrrrrrrrrr@{}}
%DIFDELCMD < \toprule
%DIFDELCMD <  & \multicolumn{4}{c}{New variant} & \multicolumn{4}{c}{New setting} & \multicolumn{2}{c}{New family} \\
%DIFDELCMD < \cmidrule(lr){2-5}\cmidrule(lr){6-9}\cmidrule(lr){10-11}
%DIFDELCMD < Variant & M1s & M2 & M5 & NN & M2 & M5 & NN & M1sn & M2 & M5 \\
%DIFDELCMD < \midrule
%DIFDELCMD < Reference & 4.9 & 23/0.15 & 7.2/0.75 & 3.4 & 18/8.1 & 19/6.8 & 8.5 & 8.2 & 51/16 & 38/28 \\
%DIFDELCMD < Temperature-adjusted & 5.0 & 23/0.10 & 7.2/0.65 & 3.4 & 17/7.6 & 19/6.8 & 9.4 & 8.9 & 48/15 & 37/27 \\
%DIFDELCMD < Warm-up dropped & 6.5 & 27/0.20 & 9.0/0.45 & 4.1 & 21/8.7 & 24/8.7 & 12 & 8.8 & 63/18 & 39/29 \\
%DIFDELCMD < Floor between series & 2.9 & 11/0 & 1.3/0.20 & 0.75 & 9.9/2.6 & 11/2.3 & 2.7 & 4.5 & 16/3.8 & 12/9.3 \\
%DIFDELCMD < Floor: batches of 25 & 4.6 & 21/0.15 & 6.4/0.70 & 2.9 & 16/6.9 & 17/6.3 & 7.7 & 7.3 & 43/14 & 31/23 \\
%DIFDELCMD < Floor: batches of 100 & 5.3 & 24/0.15 & 8.2/0.80 & 3.5 & 19/8.5 & 21/7.9 & 8.9 & 8.9 & 57/17 & 46/29 \\
%DIFDELCMD < Floor: 90\,\% quantile & 6.2 & 31/0.20 & 8.8/0.90 & 4.0 & 21/10 & 23/8.6 & 10 & 10 & 61/20 & 45/34 \\
%DIFDELCMD < Floor: 99\,\% quantile & 3.4 & 14/0.10 & 5.2/0.50 & 2.4 & 11/4.4 & 13/4.4 & 5.7 & 6.1 & 37/11 & 29/15 \\
%DIFDELCMD < Calibration: 50 parts & 5.0 & 22/0.20 & 6.8/0.70 & 3.4 & 18/7.5 & 19/6.2 & 8.4 & 7.7 & 50/16 & 37/28 \\
%DIFDELCMD < Calibration: 100 parts & 4.9 & 22/0.40 & 6.8/0.40 & 3.3 & 18/7.5 & 19/6.2 & 8.7 & 7.8 & 50/16 & 38/28 \\
%DIFDELCMD < Calibration: 250 parts & 4.9 & 22/0.45 & 7.0/0.45 & 3.3 & 18/7.5 & 19/6.2 & 8.7 & 8.3 & 50/16 & 38/28 \\
%DIFDELCMD < Conformance $p=0.90$ & 4.9 & 23/0.20 & 7.2/0.70 & 3.3 & 18/7.5 & 18/6.4 & 8.4 & 8.1 & 51/15 & 38/28 \\
%DIFDELCMD < Conformance $p=0.99$ & 5.1 & 23/0 & 7.5/0.70 & 3.3 & 18/8.3 & 20/6.9 & 8.7 & 8.4 & 51/16 & 38/28 \\
%DIFDELCMD < Interval level 0.80 &  & 23/0.15 & 6.4/1.0 &  & 18/8.1 & 17/8.1 &  &  & 50/17 & 38/28 \\
%DIFDELCMD < Interval level 0.95 &  & 23/0.15 & 8.0/0.25 &  & 18/8.1 & 20/6.2 &  &  & 51/16 & 39/27 \\
%DIFDELCMD < Response surface & 4.8 & 25/0.30 & 7.2/0.70 & 3.4 & 18/7.6 & 20/6.9 & 8.5 & 8.2 & 49/25 & 37/27 \\
%DIFDELCMD < GP: sampling bound &  &  & 6.7/1.4 &  &  & 18/7.4 &  &  &  & 38/30 \\
%DIFDELCMD < Unit: $q_{95}$ floor & 5.3 & 17/0.15 & 6.9/0.40 & 2.8 & 15/6.0 & 18/6.2 & 6.9 & 7.1 & 41/13 & 34/30 \\
%DIFDELCMD < \midrule Refit, $\delta_\mathrm{d}$ & 4.9--7.3 & 22--23 & 4.0--11 & 2.3--4.6 & 16--18 & 18--20 & 6.9--9.2 & 7.3--10 & 49--53 & 34--39 \\
%DIFDELCMD < \botrule
%DIFDELCMD < \end{tabular}
%DIFDELCMD < %%%
\DIFdelendFL \DIFaddbeginFL \begin{tabular}{@{}lrrrrrrrrrr@{}}
\toprule
 & \multicolumn{4}{c}{New variant} & \multicolumn{4}{c}{New setting} & \multicolumn{2}{c}{New family} \\
\cmidrule(lr){2-5}\cmidrule(lr){6-9}\cmidrule(lr){10-11}
Variant & M1s & M2 & M5 & NN & M2 & M5 & NN & M1sn & M2 & M5 \\
\midrule
Reference & 4.9 & 23/0.15 & 7.2/0.75 & 3.4 & 18/8.1 & 19/6.8 & 8.5 & 8.2 & 51/16 & 38/28 \\
Temperature-adjusted & 5.0 & 23/0.10 & 7.2/0.65 & 3.4 & 17/7.6 & 19/6.8 & 9.4 & 8.9 & 48/15 & 37/27 \\
Warm-up dropped & 6.5 & 27/0.20 & 9.0/0.45 & 4.1 & 21/8.7 & 24/8.7 & 12 & 8.8 & 63/18 & 39/29 \\
Floor between series & 2.9 & 11/0 & 1.3/0.20 & 0.75 & 9.9/2.6 & 11/2.3 & 2.7 & 4.5 & 16/3.8 & 12/9.3 \\
Floor: batches of 25 & 4.6 & 21/0.15 & 6.4/0.70 & 2.9 & 16/6.9 & 17/6.3 & 7.7 & 7.3 & 43/14 & 31/23 \\
Floor: batches of 100 & 5.3 & 24/0.15 & 8.2/0.80 & 3.5 & 19/8.5 & 21/7.9 & 8.9 & 8.9 & 57/17 & 46/29 \\
Floor: 90\,\% quantile & 6.2 & 31/0.20 & 8.8/0.90 & 4.0 & 21/10 & 23/8.6 & 10 & 10 & 61/20 & 45/34 \\
Floor: 99\,\% quantile & 3.4 & 14/0.10 & 5.2/0.50 & 2.4 & 11/4.4 & 13/4.4 & 5.7 & 6.1 & 37/11 & 29/15 \\
Calibration: 50 parts & 5.0 & 22/0.20 & 6.8/0.70 & 3.4 & 18/7.5 & 19/6.2 & 8.4 & 7.7 & 50/16 & 37/28 \\
Calibration: 100 parts & 4.9 & 22/0.40 & 6.8/0.40 & 3.3 & 18/7.5 & 19/6.2 & 8.7 & 7.8 & 50/16 & 38/28 \\
Calibration: 250 parts & 4.9 & 22/0.45 & 7.0/0.45 & 3.3 & 18/7.5 & 19/6.2 & 8.7 & 8.3 & 50/16 & 38/28 \\
Conformance $p=0.90$ & 4.9 & 23/0.20 & 7.2/0.70 & 3.3 & 18/7.5 & 18/6.4 & 8.4 & 8.1 & 51/15 & 38/28 \\
Conformance $p=0.99$ & 5.1 & 23/0 & 7.5/0.70 & 3.3 & 18/8.3 & 20/6.9 & 8.7 & 8.4 & 51/16 & 38/28 \\
Interval level 0.80 &  & 23/0.15 & 6.4/1.0 &  & 18/8.1 & 17/8.1 &  &  & 50/17 & 38/28 \\
Interval level 0.95 &  & 23/0.15 & 8.0/0.25 &  & 18/8.1 & 20/6.2 &  &  & 51/16 & 39/27 \\
Response surface & 4.8 & 25/0.30 & 7.2/0.70 & 3.4 & 18/7.6 & 20/6.9 & 8.5 & 8.2 & 49/25 & 37/27 \\
GP: sampling bound &  &  & 6.7/1.4 &  &  & 18/7.4 &  &  &  & 38/30 \\
Unit: $q_{95}$ floor & 5.3 & 17/0.15 & 6.9/0.40 & 2.8 & 15/6.0 & 18/6.2 & 6.9 & 7.1 & 41/13 & 34/30 \\
\midrule Refit, $\delta_\mathrm{d}$ & 4.9--7.3 & 22--23 & 4.0--11 & 2.3--4.6 & 16--18 & 18--20 & 6.9--9.2 & 7.3--10 & 49--53 & 34--39 \\
Refit, $\delta_\mathrm{s}$ & 4.9--7.3 & 0.15--1.4 & 0.75--3.0 & 2.3--4.6 & 6.8--8.4 & 6.2--7.8 & 6.9--9.2 & 7.3--10 & 15--22 & 27--33 \\
\botrule
\end{tabular}
\DIFaddendFL \footnotetext{Rules refitted for every variant; empty cells: the variant does not affect the rule; a new family in its own floor \DIFaddbeginFL \DIFaddFL{(in the produced family's floor in Table~\ref{tab:distances})}\DIFaddendFL . The floor between series contains the lubrication effect and, for a new variant, the series that NN averages, which makes it close to circular there. Warm-up: first 150 parts. Calibration: calibration $q_{95}$ from the first $n$ parts, truths and floors from the full series. Response surface: simulations from a quadratic fit over thickness and friction; other GP kernels and descriptors move M5 by at most 0.4 floors. Unit: the $q_{95}$ floor, 95\,\% quantile of the difference between the $q_{95}$ of two batches of 100 parts. Refit: 95\,\% \DIFdelbeginFL \DIFdelFL{interval }\DIFdelendFL \DIFaddbeginFL \DIFaddFL{intervals }\DIFaddendFL of the decisive \DIFdelbeginFL \DIFdelFL{distance }\DIFdelendFL \DIFaddbeginFL \DIFaddFL{and safe distances }\DIFaddendFL over 200 draws of each geometry's lubrication patterns with replacement, every rule refitted (Section~\ref{sec:evaldesign}).}
\end{table}
%% END GENERATED robust

%% BEGIN GENERATED intervals
\begin{table}[!ht]
\tabsize
\setlength{\tabcolsep}{2.6pt}
\caption{Interval rules by relation: coverage, half-width and the interval centre as a point rule}\label{tab:intervals}
\begin{tabular}{@{}lrrrrrrrrrrrr@{}}
\toprule
 & \multicolumn{3}{c}{Sibling produced} & \multicolumn{3}{c}{Interpolated setting} & \multicolumn{3}{c}{Extrapolated setting} & \multicolumn{3}{c}{New family} \\
\cmidrule(lr){2-4}\cmidrule(lr){5-7}\cmidrule(lr){8-10}\cmidrule(lr){11-13}
Rule & Cov & HW & Ctr & Cov & HW & Ctr & Cov & HW & Ctr & Cov & HW & Ctr \\
\midrule
M2 & 89 & 7.6 & 8.6 & 71 & 5.7 & 10 & 44 & 4.8 & 13 & 47 & 7.6 & 34 \\
M2n & 90 & 5.9 & 9.1 & 93 & 5.0 & 4.8 & 33 & 3.1 & 12 & 12 & 5.7 & 131 \\
M3 & 96 & 9.3 & 9.9 & 43 & 7.0 & 19 & 43 & 8.9 & 21 & 40 & 9.7 & 29 \\
M3n & 89 & 3.0 & 3.7 & 36 & 3.4 & 11 & 48 & 2.3 & 9.9 & 7 & 3.3 & 130 \\
M5 & 86 & 2.2 & 3.3 & 71 & 6.4 & 9.8 & 46 & 4.6 & 13 & 12 & 2.3 & 34 \\
M5n & 85 & 2.3 & 3.1 & 93 & 5.8 & 4.7 & 36 & 3.7 & 12 & 0 & 2.2 & 130 \\
\botrule
\end{tabular}
\footnotetext{Cov: share of cases whose true $q_{95}$ lies in the interval (\%; nominal 90). HW: median half-width (floors). Ctr: decisive distance of the interval centre used as a point rule (floors). A new family in the floor of the produced family. The GP noise variance sits at its lower bound in 97 and 99\,\% of the M5 and M5n fits with a sibling and in every new-setting and new-family fit.}
\end{table}
%% END GENERATED intervals

%% BEGIN GENERATED step6
\begin{table}[!ht]
\tabsize
\setlength{\tabcolsep}{2.2pt}
\caption{The procedure's decision rule (step 6) by relation: guard band, distances and trial shares}\label{tab:step6}
\begin{tabular}{@{}lrrrrrrrrrrrr@{}}
\toprule
 & \multicolumn{4}{c}{Sibling produced} & \multicolumn{4}{c}{Interpolated setting} & \multicolumn{4}{c}{Extrapolated setting} \\
\cmidrule(lr){2-5}\cmidrule(lr){6-9}\cmidrule(lr){10-13}
Rule & $g$ & $\delta_\mathrm{d}/\delta_\mathrm{s}$ & $T_{10}$ & $T_{20}$ & $g$ & $\delta_\mathrm{d}/\delta_\mathrm{s}$ & $T_{10}$ & $T_{20}$ & $g$ & $\delta_\mathrm{d}/\delta_\mathrm{s}$ & $T_{10}$ & $T_{20}$ \\
\midrule
M1 & none & -- & 100 & 100 & 8.75 & 20/2.1 & 44 & 4 & none & -- & 100 & 100 \\
M1s & 2.0 & 6.9/2.9 & 1 & 0 & none & -- & 100 & 100 & none & -- & 100 & 100 \\
M2 & 0 & 23/0.15 & 33 & 8 & 0.75 & 19/2.0 & 31 & 4 & none & -- & 100 & 100 \\
M5 & 0 & 7.2/0.75 & 1 & 0 & 0.5 & 20/2.3 & 40 & 4 & none & -- & 100 & 100 \\
M5n & 0 & 7.2/0.55 & 0 & 0 & 0 & 16/0 & 25 & 0 & 6.75 & 25/1.9 & 48 & 15 \\
NN & 0.25 & 3.6/3.1 & 0 & 0 & 2.0 & 6.7/2.7 & 0 & 0 & 7.0 & 17/2.5 & 26 & 1 \\
M3 & 0 & 23/0 & 45 & 6 & 17.75 & 47/0 & 94 & 69 & 16.75 & 51/0.05 & 85 & 69 \\
M3n & 0 & 11/0.05 & 9 & 0 & 6.0 & 21/1.0 & 55 & 6 & 6.5 & 19/1.3 & 47 & 3 \\
M1sn & 0.75 & 4.3/2.8 & 0 & 0 & 2.0 & 6.7/2.7 & 0 & 0 & 6.0 & 15/2.9 & 21 & 0 \\
\botrule
\end{tabular}
\footnotetext{$g$: guard band (floors) under the reference prior (Section~\ref{sec:procedure}); \DIFaddbeginFL \DIFaddFL{no applicability guard is applied; }\DIFaddendFL $\delta_\mathrm{d}/\delta_\mathrm{s}$: distances of the rule with its interval widened by $g$; $T_{10}$, $T_{20}$: share of verdicts sent to trial (\%) at requirements 10 and 20 floors from the truth. None: no band up to 20 floors reaches 95\,\%, so the procedure sends every design to trial; so also M0 and M0w in every relation, Mc with a sibling and for an extrapolated setting (for an interpolated one $g=18$), and every rule for a new family, scored in the produced family's floor.}
\end{table}
%% END GENERATED step6

\begin{table}[!ht]
\footnotesize
\setlength{\tabcolsep}{3pt}
\caption{Where the numbers quoted in the text come from: scripts and result files (\texttt{results/*.csv})}\label{tab:mapping}
\DIFdelbeginFL %DIFDELCMD < \begin{tabular}{@{}L{1.0cm}L{4.1cm}L{7.2cm}@{}}
%DIFDELCMD < \toprule
%DIFDELCMD < Section & Results & Script: result files \\
%DIFDELCMD < \midrule
%DIFDELCMD < 3.4, 4.1, 4.3 & failing-side weight, run metadata and oil film, scan-line scatter, measurement resolution & panel.py: panel\_failing\_weight, panel\_runs, panel\_resolution\_steps; measurement.py: meas\_scan \\
%DIFDELCMD < 5.1 & time order, lag profile, normal model, subgroups, part-to-part share, between series, $q_{95}$ reproducibility, in-line drift & alternatives.py: alt\_floor; sensitivity.py: sens\_drift; panel.py: panel\_variogram, panel\_floor\_protocol; measurement.py: meas\_floor; robustness.py: rob\_floor\_between; inline.py: inline\_drift \\
%DIFDELCMD < 5.2 & resolvable contrasts and changes, simulation offsets, force effect, arm sign, cup depth & alternatives.py: alt\_effects, alt\_mrc; robustness.py: rob\_bhf\_pairs; decisions.py: dec\_alternatives; panel.py: panel\_force\_effect, panel\_arm\_sign, panel\_m0\_split \\
%DIFDELCMD < 5.3--5.4 & distances in both floors, paired differences, strata, force levels, decomposition, pooled coverage & decisions.py: dec\_resolution, dec\_paired, dec\_coverage; panel.py: panel\_source\_floor(\_paired), panel\_sibling\_strata, panel\_force\_levels, panel\_decomposition \\
%DIFDELCMD < 5.5 & simulation pairs, drawing nominal, surrogate, coverage, GP fits & decisions.py: dec\_paired, dec\_coverage, dec\_gp; panel.py: panel\_nominal\_offset, panel\_source\_floor(\_paired); design\_space.py: ds\_surrogate(\_rule) \\
%DIFDELCMD < 5.6, 5.9 & budget, short series, robustness variants, refit, $q_{95}$ unit & budget.py: budget\_design, budget\_resolution; robustness.py: rob\_decisions, rob\_refit(\_paired); panel.py: panel\_q95\_unit \\
%DIFDELCMD < 5.7 & margins, guard bands and priors, guards, worked example & panel.py: panel\_step6(\_curves), panel\_guard\_priors; guard.py: dec\_guard; decisions.py: dec\_intervals \\
%DIFDELCMD < 5.8 & capability-anchored limits, ISO 2768-1 scenario, costs & panel.py: panel\_capability; scenarios.py: scen\_strata, scen\_scores, scen\_cost\_map \\
%DIFDELCMD < 6.5 & PA12 attempt & case2\_pa12.py: case2\_floor, case2\_effects \\
%DIFDELCMD < \botrule
%DIFDELCMD < \end{tabular}
%DIFDELCMD < %%%
\DIFdelendFL \DIFaddbeginFL \begin{tabular}{@{}L{1.0cm}L{4.1cm}L{7.2cm}@{}}
\toprule
Section & Results & Script: result files \\
\midrule
3.4, 4.1, 4.3 & failing-side weight, run metadata and oil film, scan-line scatter, measurement resolution & panel.py: panel\_failing\_weight, panel\_runs, panel\_resolution\_steps; measurement.py: meas\_scan \\
5.1 & time order, lag profile, normal model, subgroups, part-to-part share, between series, $q_{95}$ reproducibility, in-line drift & alternatives.py: alt\_floor; sensitivity.py: sens\_drift; panel.py: panel\_variogram, panel\_floor\_protocol; measurement.py: meas\_floor; robustness.py: rob\_floor\_between; inline.py: inline\_drift \\
5.2 & resolvable contrasts and changes, simulation offsets, force effect, arm sign, cup depth & alternatives.py: alt\_effects, alt\_mrc; robustness.py: rob\_bhf\_pairs; decisions.py: dec\_alternatives; panel.py: panel\_force\_effect, panel\_arm\_sign, panel\_m0\_split \\
5.3--5.4 & distances in both floors, paired differences, strata, force levels, decomposition, pooled coverage & decisions.py: dec\_resolution, dec\_paired, dec\_coverage; panel.py: panel\_source\_floor(\_paired), panel\_sibling\_strata, panel\_force\_levels, panel\_decomposition \\
5.5 & simulation pairs, drawing nominal, surrogate, coverage, GP fits & decisions.py: dec\_paired, dec\_coverage, dec\_gp; panel.py: panel\_nominal\_offset, panel\_source\_floor(\_paired); design\_space.py: ds\_surrogate(\_rule) \\
5.6, 5.9 & budget, short series, robustness variants, refit, $q_{95}$ unit & budget.py: budget\_design, budget\_resolution; robustness.py: rob\_decisions, rob\_refit(\_paired); panel.py: panel\_q95\_unit \\
3.6, 5.7 & margins, guard bands and priors, uncapped bands, guards, worked example & panel.py: panel\_step6(\_curves), panel\_guard\_priors; guard.py: dec\_guard; decisions.py: dec\_intervals \\
5.8 & capability-anchored limits, ISO 2768-1 scenario, costs & panel.py: panel\_capability; scenarios.py: scen\_strata, scen\_scores, scen\_cost\_map \\
6.5 & PA12 attempt & case2\_pa12.py: case2\_floor, case2\_effects \\
\botrule
\end{tabular}
\DIFaddendFL \footnotetext{\texttt{scripts/run\_all.sh} runs the scripts in order; \texttt{make\_tables.py} writes every table and \texttt{figs.py} every figure from the result files.}
\end{table}

\end{appendices}

\backmatter

\bmhead{Acknowledgements}
The author thanks the authors of the RDDAC and DDACS datasets at the University of Stuttgart for publishing the data on which this work rests.

\section*{Declarations}

\bmhead{Funding}
No specific funding was received for this work.

\bmhead{Competing interests}
The author declares no competing interests.

\bmhead{Ethics approval and consent}
Not applicable.

\bmhead{Data availability}
All data are public. RDDAC: doi:10.18419/DARUS-5589 (version 2.0); DDACS: doi:10.18419/DARUS-4801 (version 3.0); both licensed CC BY 4.0. The extracted per-part and per-simulation feature tables and all result files are archived with the code \citep{canseven2026archive}.

\bmhead{Code availability}
The complete analysis pipeline, which reproduces every number in this paper from the public data, is archived with the feature tables and result files of this version on Zenodo \citep{canseven2026archive}; Table~\ref{tab:mapping} maps the numbers to its scripts and files. A working copy is at \url{https://github.com/htcanseven/Claude} (directory \texttt{scripts/}).

\bmhead{Author contributions}
H.T.C. conceived the method, designed and carried out the analysis, and wrote the paper.

\bmhead{Use of generative AI}
A large language model (Claude, Anthropic) was used under the author's direction for code development, literature retrieval, data analysis, drafting and editing, and simulated peer review of earlier versions. The author designed the study, verified all analyses, checked every reference against its published record, and takes full responsibility for the content. All figures are produced by the released code from the data; no AI-generated images are used.


%% ----------------------------------------------------------------------
\FloatBarrier
{\small
\bibliography{refs}}

\end{document}
